\documentclass[11pt,a4paper]{article}

\usepackage[margin=1in]{geometry}
\usepackage[T1]{fontenc}
\usepackage[utf8]{inputenc}
\usepackage{lmodern}
\usepackage{microtype}
\usepackage{amsmath,amssymb,amsfonts,amsthm,mathtools,bm}
\usepackage{enumitem}
\usepackage{hyperref}
\usepackage{booktabs}
\usepackage{graphicx}
\usepackage{xcolor}

\hypersetup{
    colorlinks=true,
    linkcolor=blue,
    citecolor=blue,
    urlcolor=blue,
    pdftitle={Convergence Analysis of HRAC},
    pdfauthor={OpenAI}
}

\title{\textbf{Convergence Analysis of HRAC}\\
\large A Standalone \LaTeX{} File for Direct Overleaf Compilation}
\author{}
\date{\today}

\newtheorem{assumption}{Assumption}
\newtheorem{lemma}{Lemma}
\newtheorem{theorem}{Theorem}
\newtheorem{corollary}{Corollary}
\newtheorem{proposition}{Proposition}
\newtheorem{definition}{Definition}
\newtheorem{condition}{Condition}
\theoremstyle{remark}
\newtheorem{remark}{Remark}

\newcommand{\R}{\mathbb{R}}
\newcommand{\E}{\mathbb{E}}
\newcommand{\F}{\mathcal{F}}
\newcommand{\HH}{\mathcal{H}}
\newcommand{\BB}{\mathcal{B}}
\newcommand{\clip}{\operatorname{clip}}
\newcommand{\med}{\operatorname{med}}
\newcommand{\norm}[1]{\left\lVert #1 \right\rVert}
\newcommand{\abs}[1]{\left\lvert #1 \right\rvert}

\begin{document}

\maketitle

\begin{abstract}
This note gives a modular convergence analysis for \emph{History-Residual
Adaptive Clipping} (HRAC) in Byzantine-robust federated optimisation. The
optimization layer compares the HRAC aggregate with the honest average and
derives smooth non-convex descent with an explicit aggregation-error floor.
Honest-majority median protection and bounded honest updates imply
deterministic boundedness of the clipped messages and of \(\E\norm{g_t}^2\).
For the adaptive statistics, the main practically relevant conclusion is
late-stage Byzantine group-level \(\nu\)-band elevation, which induces stronger
penalty pressure and smaller weights for Byzantine clients once the activation
threshold is reached. In the experiments motivating this note, neither honest
nor Byzantine \(d\)- or \(\nu\)-statistics converge to \(0\); instead, both
groups remain in positive late-stage bands, with Byzantine clients typically
occupying a higher \(\nu\)-band. Under additional spread control, this
group-level band elevation upgrades to pointwise \(\nu\)-separation. A windowed
\(d\)-gap certificate further turns logged statistics into an exact sufficient
condition for blockwise \(\nu\)-separation. The file is standalone and
directly compilable in Overleaf.
\end{abstract}

\tableofcontents

\section{Introduction}

The goal of this document is to give a mathematically clean convergence
analysis for HRAC. Rather than proving the code line by line, we abstract HRAC
as a robust aggregation rule and proceed in two steps:

\begin{enumerate}[label=\arabic*.]
    \item show that the HRAC aggregate remains close to the ideal honest
    average update;
    \item show that such an aggregation error can be absorbed into the descent
    analysis of smooth non-convex optimisation.
\end{enumerate}

For the adaptive statistics, the main theorem-level target is not universal
pointwise ordering between Byzantine and honest clients at every late-stage
round. Instead, the primary claim is the weaker and empirically better-matched
one: Byzantine clients exhibit a larger late-stage group-mean \(\nu\)-band,
which induces stronger penalty pressure and smaller weights once the activation
threshold becomes active. In particular, the intended mechanism is not that
honest statistics vanish while Byzantine statistics remain nonzero. Rather,
both groups may stay in positive late-stage bands, while Byzantine clients are
typically elevated relative to the honest population. Pointwise separation is
kept as a stronger sufficient refinement.

\section{Problem Setup}

Let \(x_t \in \R^d\) denote the global model at round \(t\).
Suppose there are \(N\) participating clients, partitioned into an honest set
\(\HH\) and a Byzantine set \(\BB\), with
\[
    |\HH| = H, \qquad |\BB| = B, \qquad N = H + B.
\]

At round \(t\), the server receives client updates
\[
    \{\Delta_i^t\}_{i=1}^N \subset \R^d,
\]
and performs the global update
\begin{equation}
    x_{t+1} = x_t - \eta_t g_t,
    \label{eq:server_update}
\end{equation}
where \(g_t\) is the aggregate produced by HRAC.

The ideal honest average update is defined as
\begin{equation}
    \bar{\Delta}_{\HH}^t
    :=
    \frac{1}{H} \sum_{i \in \HH} \Delta_i^t.
    \label{eq:honest_average}
\end{equation}

The objective is to minimise a possibly non-convex function
\[
    f : \R^d \to \R.
\]

\section{HRAC Aggregation Rule}

HRAC operates in three stages.

\subsection{Global norm clipping}

Define the round-wise cap
\begin{equation}
    B_t := c_g \cdot \med_{1 \le j \le N} \norm{\Delta_j^t},
    \label{eq:global_cap}
\end{equation}
and the globally clipped update
\begin{equation}
    \hat{\Delta}_i^t
    :=
    \clip(\Delta_i^t, B_t)
    :=
    \begin{cases}
        \Delta_i^t, & \norm{\Delta_i^t} \le B_t,\\[2mm]
        B_t \dfrac{\Delta_i^t}{\norm{\Delta_i^t}}, & \norm{\Delta_i^t} > B_t.
    \end{cases}
    \label{eq:global_clip}
\end{equation}

\subsection{Per-client history-residual clipping}

For each client \(i\), HRAC maintains a history state
\[
    (b_i^t,\mu_i^t,\nu_i^t),
\]
where \(b_i^t\) is a running centre, \(\mu_i^t\) is a scale estimate, and
\(\nu_i^t\) is a change statistic.

Define the residual
\begin{equation}
    r_i^t := \hat{\Delta}_i^t - b_i^t,
    \label{eq:residual}
\end{equation}
and the adaptive residual clipping threshold
\begin{equation}
    \tau_i^t := c \mu_i^t.
    \label{eq:tau}
\end{equation}

The clipped residual is
\begin{equation}
    \bar{r}_i^t
    :=
    \clip(r_i^t, \tau_i^t)
    =
    \begin{cases}
        r_i^t, & \norm{r_i^t} \le \tau_i^t,\\[2mm]
        \tau_i^t \dfrac{r_i^t}{\norm{r_i^t}}, & \norm{r_i^t} > \tau_i^t.
    \end{cases}
    \label{eq:residual_clip}
\end{equation}

The post-processed message becomes
\begin{equation}
    \tilde{\Delta}_i^t := b_i^t + \bar{r}_i^t.
    \label{eq:delta_tilde}
\end{equation}

The history centre \(b_i^t\) follows the implemented EMA dynamics:
for a newly appearing client, \(b_i^t\) is initialized by
\(\hat{\Delta}_i^t\); for later rounds,
\begin{equation}
    b_i^{t+1}
    =
    \rho_b\, b_i^t + (1-\rho_b)\,\Delta_{i,\mathrm{safe}}^t,
    \label{eq:b_update}
\end{equation}
where \(\Delta_{i,\mathrm{safe}}^t\) is the aggregation-path message used in
the server history update.

% In the current implementation used in experiments
% (\texttt{enable\_post\_residual\_b\_cap=False}), no additional post-cap is
% applied to \(\tilde{\Delta}_i^t\) after residual reconstruction.

\subsection{Adaptive weighting}

HRAC assigns weights \(w_i^t \ge 0\) satisfying
\begin{equation}
    \sum_{i=1}^N w_i^t = 1,
    \label{eq:weight_simplex}
\end{equation}
where the weights are determined from the statistics
\(\{\nu_i^t\}_{i=1}^N\).

The final aggregate is
\begin{equation}
    g_t = \sum_{i=1}^N w_i^t \tilde{\Delta}_i^t.
    \label{eq:aggregate}
\end{equation}

\section{Assumptions}

\begin{assumption}[Lower bounded objective]
\label{ass:lower_bounded}
\(f\) is bounded below: \(f(x)\ge f^\star\) for all \(x\in\R^d\).
\end{assumption}

\begin{assumption}[\(L\)-smoothness]
\label{ass:smooth}
\(\norm{\nabla f(x)-\nabla f(y)}\le L\norm{x-y}\) for all \(x,y\in\R^d\).
\end{assumption}

\begin{assumption}[Bias and variance of the honest average]
\label{ass:bias_variance}
Conditioned on \(\F_t\), the honest average has bounded bias and variance:
\begin{equation}
    \norm{
        \E\!\left[\bar{\Delta}_{\HH}^t \mid \F_t\right] - \nabla f(x_t)
    }
    \le \xi,
    \label{eq:bias}
\end{equation}
\begin{equation}
    \E\!\left[
        \norm{
            \bar{\Delta}_{\HH}^t -
            \E\!\left[\bar{\Delta}_{\HH}^t \mid \F_t\right]
        }^2
        \,\middle|\, \F_t
    \right]
    \le \frac{\sigma^2}{H}.
    \label{eq:variance}
\end{equation}
\end{assumption}

\begin{assumption}[Bounded honest update magnitude]
\label{ass:honest_update_bound}
There exists \(G_H<\infty\) such that for all honest clients \(i\in\HH\) and
all rounds \(t\),
\begin{equation}
    \norm{\Delta_i^t}\le G_H.
    \label{eq:honest_update_bound}
\end{equation}
\end{assumption}

\begin{assumption}[Honest majority for robust median cap]
\label{ass:honest_majority}
The Byzantine population is strictly less than the honest population:
\begin{equation}
    B < H
    \qquad(\text{equivalently } B < N/2).
    \label{eq:honest_majority}
\end{equation}
\end{assumption}

\section{Basic Properties of HRAC}

\begin{lemma}[Boundedness of globally clipped inputs]
\label{lem:bounded}
For every round \(t\) and every client \(i\),
\begin{equation}
    \norm{\hat{\Delta}_i^t} \le B_t,
    \label{eq:uniform_bound}
\end{equation}
\end{lemma}

\begin{proof}
The bound \(\norm{\hat{\Delta}_i^t} \le B_t\) follows directly from the
definition of global clipping in~\eqref{eq:global_clip}.

\end{proof}

\begin{lemma}[Median-cap protection under honest majority]
\label{lem:median_cap_honest}
Under Assumption~\ref{ass:honest_majority}, for every round \(t\), the global
median norm in \eqref{eq:global_cap} is bounded between two honest norms.
Consequently,
\begin{equation}
    B_t
    =
    c_g\,\med_{1\le j\le N}\norm{\Delta_j^t}
    \le
    c_g\,\max_{i\in\HH}\norm{\Delta_i^t}.
    \label{eq:median_honest_bound}
\end{equation}
\end{lemma}

\begin{proof}
Because \(H>B\), more than half of the \(N\) samples are honest, so every
median order statistic is bracketed by honest samples. In particular, the
median cannot exceed the maximum honest norm, giving
\eqref{eq:median_honest_bound}.
\end{proof}

\begin{corollary}[Uniform cap bound from honest majority and bounded honest updates]
\label{cor:Bmax_from_honest}
Under Assumptions~\ref{ass:honest_majority} and
\ref{ass:honest_update_bound},
\begin{equation}
    B_t
    \le
    c_g G_H
    =: B_{\max},
    \qquad \forall t.
    \label{eq:Bmax_from_honest}
\end{equation}
\end{corollary}

\begin{proof}
By Lemma~\ref{lem:median_cap_honest},
\(
B_t\le c_g\max_{i\in\HH}\norm{\Delta_i^t}
\).
Apply \eqref{eq:honest_update_bound}.
\end{proof}

\begin{lemma}[EMA boundedness for \(\mu_i^t\)]
\label{lem:mu_ema_bounded}
Suppose nonnegative scalars \(s_t\) satisfy \(s_t\le S_{\max}\) for all \(t\),
and \(\mu_{t+1}=\rho_\mu\mu_t+(1-\rho_\mu)s_t\) with \(0<\rho_\mu<1\) and
\(\mu_0\le S_{\max}\). Then \(\mu_t\le S_{\max}\) for all \(t\).
\end{lemma}

\begin{proof}
Induction with convexity:
\[
    \mu_{t+1}
    =
    \rho_\mu\mu_t+(1-\rho_\mu)s_t
    \le
    \rho_\mu S_{\max}+(1-\rho_\mu)S_{\max}
    =
    S_{\max}.
\]
\end{proof}

\begin{lemma}[Bounded history states under bounded global cap]
\label{lem:mu_tau_bounded}
Assume there exists \(B_{\max}<\infty\) such that \(B_t\le B_{\max}\) for all
rounds \(t\), and the server-side history update uses the capped safe message
(\(\norm{\Delta_{i,\mathrm{safe}}^t}\le B_t\), as in the implementation).
Then for every client \(i\) and every \(t\):
\begin{equation}
    \norm{b_i^t}\le B_{\max},
    \qquad
    \norm{r_i^t}\le 2B_{\max},
    \qquad
    \mu_i^t\le 2B_{\max},
    \qquad
    \tau_i^t=c\mu_i^t\le 2cB_{\max}.
    \label{eq:mu_tau_bounds}
\end{equation}
Moreover, the post-processed message is deterministically bounded:
\begin{equation}
    \norm{\tilde\Delta_i^t}\le (1+2c)B_{\max} := M_{\tilde\Delta}.
    \label{eq:tilde_delta_det_bound}
\end{equation}
\end{lemma}

\begin{proof}
For \(b_i^t\), the initialization satisfies \(\norm{b_i^0}\le B_{\max}\) since
it is set from a globally clipped message. If \(\norm{b_i^t}\le B_{\max}\), then
\[
    b_i^{t+1}
    =
    \rho_b b_i^t + (1-\rho_b)\Delta_{i,\mathrm{safe}}^t,
    \qquad
    \norm{\Delta_{i,\mathrm{safe}}^t}\le B_t\le B_{\max},
\]
so convexity gives \(\norm{b_i^{t+1}}\le B_{\max}\). Hence
\(\norm{b_i^t}\le B_{\max}\) for all \(t\).

Because \(\norm{\hat\Delta_i^t}\le B_t\le B_{\max}\),
\[
    \norm{r_i^t}
    =
    \norm{\hat\Delta_i^t-b_i^t}
    \le
    \norm{\hat\Delta_i^t}+\norm{b_i^t}
    \le
    2B_{\max}.
\]
If \(\mu_i^t\) is initialized at most \(2B_{\max}\) and updated by EMA of
bounded residual magnitudes (implementation path), then
Lemma~\ref{lem:mu_ema_bounded} gives \(\mu_i^t\le 2B_{\max}\); thus
\(\tau_i^t=c\mu_i^t\le 2cB_{\max}\).

Finally, \(\tilde\Delta_i^t=b_i^t+\bar r_i^t\) with
\(\norm{\bar r_i^t}\le \tau_i^t\), so
\[
    \norm{\tilde\Delta_i^t}
    \le
    \norm{b_i^t}+\norm{\bar r_i^t}
    \le
    B_{\max}+2cB_{\max}
    =
    (1+2c)B_{\max}.
\]
\end{proof}

\begin{lemma}[Aggregate second moment from deterministic cap chain]
\label{lem:g_second_derived}
Under Assumptions~\ref{ass:honest_majority} and
\ref{ass:honest_update_bound}, for every \(t\), the aggregation output satisfies
\begin{equation}
    \E\!\left[\norm{g_t}^2 \mid \F_t\right]
    \le
    M_{\tilde\Delta}^2,
    \qquad\text{hence}\qquad
    \frac{1}{T}\sum_{t=0}^{T-1}\E\norm{g_t}^2 \le M_{\tilde\Delta}^2.
    \label{eq:Gg_derived}
\end{equation}
\end{lemma}

\begin{proof}
By Corollary~\ref{cor:Bmax_from_honest}, \(B_t\le B_{\max}\) for all \(t\). 
Lemma~\ref{lem:mu_tau_bounded} then gives \(\norm{\tilde\Delta_i^t}\le M_{\tilde\Delta}\) deterministically for every client \(i\). 
Write \(g_t=\sum_{i=1}^N w_i^t \tilde{\Delta}_i^t\) with \(\sum_i w_i^t=1\),
\(w_i^t\ge 0\). Jensen's inequality yields:
\[
    \norm{g_t}^2 = \norm{\sum_{i=1}^N w_i^t \tilde{\Delta}_i^t}^2
    \le
    \sum_{i=1}^N w_i^t \norm{\tilde{\Delta}_i^t}^2
    \le 
    \sum_{i=1}^N w_i^t M_{\tilde\Delta}^2 = M_{\tilde\Delta}^2.
\]
Taking expectations preserves the bound.
\end{proof}

\begin{lemma}[Residual clipping bias identity]
\label{lem:clip_identity}
For every client \(i\) and round \(t\),
\begin{equation}
    \norm{\bar{r}_i^t - r_i^t}
    =
    (\norm{r_i^t} - \tau_i^t)_+,
    \label{eq:clip_identity}
\end{equation}
where \((a)_+ := \max\{a,0\}\).
As a consequence,
\begin{equation}
    \norm{\tilde{\Delta}_i^t - \hat{\Delta}_i^t}
    =
    (\norm{r_i^t} - \tau_i^t)_+.
    \label{eq:delta_clip_identity}
\end{equation}
\end{lemma}

\begin{proof}
If \(\norm{r_i^t} \le \tau_i^t\), then \(\bar{r}_i^t = r_i^t\), so the
left-hand side is \(0\), equal to
\((\norm{r_i^t}-\tau_i^t)_+\).

If \(\norm{r_i^t} > \tau_i^t\), then
\[
    \bar{r}_i^t = r_i^t \cdot \frac{\tau_i^t}{\norm{r_i^t}},
\]
hence
\[
    \norm{\bar{r}_i^t - r_i^t}
    =
    \left\|
        r_i^t\left(\frac{\tau_i^t}{\norm{r_i^t}} - 1\right)
    \right\|
    =
    \norm{r_i^t} - \tau_i^t.
\]
This proves~\eqref{eq:clip_identity}.
Since
\[
    \tilde{\Delta}_i^t - \hat{\Delta}_i^t
    =
    (b_i^t+\bar{r}_i^t) - (b_i^t+r_i^t)
    =
    \bar{r}_i^t-r_i^t,
\]
we obtain~\eqref{eq:delta_clip_identity}.
\end{proof}

\section{Aggregation Error Analysis}

Define the aggregation error
\begin{equation}
    e_t := g_t - \bar{\Delta}_{\HH}^t.
    \label{eq:e_t}
\end{equation}

\begin{lemma}[Aggregation error decomposition]
\label{lem:decomp}
The error \(e_t\) admits the decomposition
\begin{equation}
    e_t
    =
    \underbrace{
    \sum_{i \in \HH} w_i^t(\tilde{\Delta}_i^t - \Delta_i^t)
    }_{=:A_t}
    +
    \underbrace{
    \sum_{i \in \HH}
    \left(w_i^t - \frac{1}{H}\right)\Delta_i^t
    }_{=:W_t}
    +
    \underbrace{
    \sum_{j \in \BB} w_j^t \tilde{\Delta}_j^t
    }_{=:M_t}.
    \label{eq:decomposition}
\end{equation}
\end{lemma}

\begin{proof}
From~\eqref{eq:aggregate},
\[
    g_t
    =
    \sum_{i \in \HH} w_i^t \tilde{\Delta}_i^t
    +
    \sum_{j \in \BB} w_j^t \tilde{\Delta}_j^t.
\]
Subtract \(\bar{\Delta}_{\HH}^t = \frac{1}{H}\sum_{i\in\HH}\Delta_i^t\) and
add/subtract \(\sum_{i\in\HH} w_i^t \Delta_i^t\) to obtain
\[
    g_t - \bar{\Delta}_{\HH}^t
    =
    \sum_{i \in \HH} w_i^t(\tilde{\Delta}_i^t-\Delta_i^t)
    +
    \sum_{i \in \HH}
    \left(w_i^t-\frac{1}{H}\right)\Delta_i^t
    +
    \sum_{j \in \BB} w_j^t \tilde{\Delta}_j^t.
\]
\end{proof}

\begin{lemma}[Bounding each term of the decomposition]
\label{lem:term_bounds}
\begin{align}
    \norm{A_t}^2
    &\le
    \sum_{i \in \HH} w_i^t \norm{\tilde{\Delta}_i^t-\Delta_i^t}^2,
    \label{eq:A_t_bound}\\
    \norm{W_t}^2
    &\le
    \varepsilon_{w,t}^2 \cdot \max_{i\in\HH}\norm{\Delta_i^t}^2,
    \label{eq:W_t_bound}\\
    \norm{M_t}^2
    &\le
    \beta_t \sum_{j\in\BB} w_j^t \norm{\tilde{\Delta}_j^t}^2.
    \label{eq:M_t_bound}
\end{align}
Consequently,
\begin{equation}
    \norm{e_t}^2
    \le
    3\norm{A_t}^2 + 3\norm{W_t}^2 + 3\norm{M_t}^2.
    \label{eq:e_abc}
\end{equation}
\end{lemma}

\begin{proof}
For \(A_t\), Jensen's inequality gives
\[
    \norm{A_t}^2
    =
    \left\|
        \sum_{i \in \HH} w_i^t(\tilde{\Delta}_i^t-\Delta_i^t)
    \right\|^2
    \le
    \sum_{i \in \HH} w_i^t \norm{\tilde{\Delta}_i^t-\Delta_i^t}^2.
\]
More explicitly, since \(\sum_{i\in\HH} w_i^t = 1-\beta_t \le 1\),
weighted Cauchy--Schwarz yields
\[
    \norm{A_t}^2
    \le
    \left(\sum_{i\in\HH} w_i^t\right)
    \left(\sum_{i\in\HH} w_i^t \norm{\tilde{\Delta}_i^t-\Delta_i^t}^2\right)
    \le
    \sum_{i\in\HH} w_i^t \norm{\tilde{\Delta}_i^t-\Delta_i^t}^2.
\]

For \(W_t\),
\[
    \norm{W_t}
    \le
    \sum_{i \in \HH}
    \abs{w_i^t-\frac{1}{H}} \norm{\Delta_i^t}
    \le
    \varepsilon_{w,t} \max_{i\in\HH}\norm{\Delta_i^t},
\]
which yields~\eqref{eq:W_t_bound} after squaring.

For \(M_t\), by Cauchy--Schwarz and \(\sum_{j\in\BB}w_j^t=\beta_t\),
\[
    \norm{M_t}
    \le
    \sum_{j\in\BB} w_j^t \norm{\tilde{\Delta}_j^t}
\]
and therefore
\[
    \norm{M_t}^2
    \le
    \left(\sum_{j\in\BB}w_j^t\right)
    \left(\sum_{j\in\BB}w_j^t\norm{\tilde{\Delta}_j^t}^2\right)
    =
    \beta_t\sum_{j\in\BB}w_j^t\norm{\tilde{\Delta}_j^t}^2.
\]
which is exactly~\eqref{eq:M_t_bound}.

Finally,
\[
    e_t = A_t + W_t + M_t,
\]
so
\[
    \norm{e_t}^2
    \le
    3\norm{A_t}^2 + 3\norm{W_t}^2 + 3\norm{M_t}^2.
\]
\end{proof}

\begin{lemma}[Explicit honest clipping bias bound]
\label{lem:honest_clip}
For every honest client \(i \in \HH\),
\begin{equation}
    \norm{\tilde{\Delta}_i^t-\Delta_i^t}^2
    \le
    2\norm{\hat{\Delta}_i^t-\Delta_i^t}^2
    +
    2(\norm{r_i^t}-\tau_i^t)_+^2.
    \label{eq:honest_clip_bound}
\end{equation}
Therefore,
\begin{equation}
    \sum_{i \in \HH} w_i^t \norm{\tilde{\Delta}_i^t-\Delta_i^t}^2
    \le
    2\sum_{i \in \HH} w_i^t \norm{\hat{\Delta}_i^t-\Delta_i^t}^2
    +
    2\sum_{i \in \HH} w_i^t (\norm{r_i^t}-\tau_i^t)_+^2.
    \label{eq:honest_clip_sum}
\end{equation}
\end{lemma}

\begin{proof}
Using
\[
    \tilde{\Delta}_i^t-\Delta_i^t
    =
    (\tilde{\Delta}_i^t-\hat{\Delta}_i^t)
    +
    (\hat{\Delta}_i^t-\Delta_i^t),
\]
we have
\[
    \norm{\tilde{\Delta}_i^t-\Delta_i^t}^2
    \le
    2\norm{\tilde{\Delta}_i^t-\hat{\Delta}_i^t}^2
    +
    2\norm{\hat{\Delta}_i^t-\Delta_i^t}^2.
\]
Lemma~\ref{lem:clip_identity} gives
\[
    \norm{\tilde{\Delta}_i^t-\hat{\Delta}_i^t}^2
    =
    (\norm{r_i^t}-\tau_i^t)_+^2.
\]
Substituting yields~\eqref{eq:honest_clip_bound}. Summing with weights
\(w_i^t\) gives~\eqref{eq:honest_clip_sum}.
\end{proof}

\begin{lemma}[Overall aggregation error bound]
\label{lem:agg_bound}
\begin{align}
    \norm{e_t}^2
    \le\;&
    6\sum_{i \in \HH} w_i^t \norm{\hat{\Delta}_i^t-\Delta_i^t}^2
    +
    6\sum_{i \in \HH} w_i^t (\norm{r_i^t}-\tau_i^t)_+^2
    \nonumber\\
    &+
    3\varepsilon_{w,t}^2 \max_{i\in\HH}\norm{\Delta_i^t}^2
    +
    3\beta_t \sum_{j\in\BB} w_j^t \norm{\tilde{\Delta}_j^t}^2.
    \label{eq:agg_error_final}
\end{align}
\end{lemma}

\begin{proof}
Combine Lemma~\ref{lem:term_bounds} with
Lemma~\ref{lem:honest_clip}.
\end{proof}

\section{From Honest Average to the True Gradient}

\begin{lemma}[Deviation of the honest average from the true gradient]
\label{lem:honest_gradient}
Under Assumption~\ref{ass:bias_variance},
\begin{equation}
    \E\!\left[
        \norm{\bar{\Delta}_{\HH}^t - \nabla f(x_t)}^2
        \,\middle|\, \F_t
    \right]
    \le
    \frac{\sigma^2}{H} + \xi^2.
    \label{eq:honest_gradient_gap}
\end{equation}
\end{lemma}

\begin{proof}
Write
\[
    \bar{\Delta}_{\HH}^t - \nabla f(x_t)
    =
    \left(
        \bar{\Delta}_{\HH}^t -
        \E[\bar{\Delta}_{\HH}^t \mid \F_t]
    \right)
    +
    \left(
        \E[\bar{\Delta}_{\HH}^t \mid \F_t] - \nabla f(x_t)
    \right).
\]
Using the variance bound~\eqref{eq:variance} and the bias bound~\eqref{eq:bias}
gives the result.
\end{proof}

\section{One-Step Descent Analysis}

\begin{lemma}[One-step descent inequality]
\label{lem:descent}
Suppose Assumptions~\ref{ass:lower_bounded} to \ref{ass:honest_majority}
hold.
Then, for any step size \(\eta_t > 0\),
\begin{align}
    \E[f(x_{t+1})]
    \le\;&
    \E[f(x_t)]
    -
    \frac{\eta_t}{2}\E\norm{\nabla f(x_t)}^2
    \nonumber\\
    &+
    \eta_t \E\norm{e_t}^2
    +
    \eta_t\left(\frac{\sigma^2}{H}+\xi^2\right)
    +
    \frac{L\eta_t^2}{2}\E\norm{g_t}^2.
    \label{eq:descent}
\end{align}
\end{lemma}

\begin{proof}
By Assumption~\ref{ass:smooth},
\[
    f(x_{t+1})
    \le
    f(x_t)
    +
    \langle \nabla f(x_t), x_{t+1}-x_t \rangle
    +
    \frac{L}{2}\norm{x_{t+1}-x_t}^2.
\]
Substituting \(x_{t+1}=x_t-\eta_t g_t\) gives
\[
    f(x_{t+1})
    \le
    f(x_t)
    -
    \eta_t \langle \nabla f(x_t), g_t \rangle
    +
    \frac{L\eta_t^2}{2}\norm{g_t}^2.
\]

Using the identity \(-\langle a, b \rangle = -\frac{1}{2}\norm{a}^2 - \frac{1}{2}\norm{b}^2 + \frac{1}{2}\norm{a-b}^2\), we can rewrite the inner product exactly as:
\[
    -\langle \nabla f(x_t), g_t \rangle
    =
    -\frac{1}{2}\norm{\nabla f(x_t)}^2
    - \frac{1}{2}\norm{g_t}^2
    + \frac{1}{2}\norm{g_t - \nabla f(x_t)}^2.
\]
Since \(-\frac{1}{2}\norm{g_t}^2 \le 0\), we drop it for an upper bound:
\[
    -\langle \nabla f(x_t), g_t \rangle
    \le
    -\frac{1}{2}\norm{\nabla f(x_t)}^2
    + \frac{1}{2}\norm{g_t - \nabla f(x_t)}^2.
\]
Now, using the decomposition \(g_t = \bar{\Delta}_{\HH}^t + e_t\), and applying the inequality \(\norm{a+b}^2 \le 2\norm{a}^2 + 2\norm{b}^2\), we bound the distance between the aggregate and the true gradient:
\[
    \norm{g_t - \nabla f(x_t)}^2
    =
    \norm{e_t + \bar{\Delta}_{\HH}^t - \nabla f(x_t)}^2
    \le
    2\norm{e_t}^2 + 2\norm{\bar{\Delta}_{\HH}^t - \nabla f(x_t)}^2.
\]
Substituting this back into the smoothness inequality yields:
\[
    f(x_{t+1})
    \le
    f(x_t)
    - \frac{\eta_t}{2}\norm{\nabla f(x_t)}^2
    + \eta_t\norm{e_t}^2
    + \eta_t\norm{\bar{\Delta}_{\HH}^t - \nabla f(x_t)}^2
    + \frac{L\eta_t^2}{2}\norm{g_t}^2.
\]
Taking the conditional expectation \(\E[\cdot \mid \F_t]\) on both sides and applying Lemma~\ref{lem:honest_gradient} for the honest average deviation yields the final result \eqref{eq:descent}.
\end{proof}

\section{Main Convergence Theorem}

\begin{theorem}[Convergence under effective-adversary conditions]
\label{thm:main}
Suppose Assumptions~\ref{ass:lower_bounded} to \ref{ass:honest_majority} hold.
Let the step size be constant and satisfy
\begin{equation}
    0 < \eta \le \frac{1}{2L}.
    \label{eq:stepsize}
\end{equation}
Then the iterates generated by HRAC satisfy
\begin{align}
    \frac{1}{T}\sum_{t=0}^{T-1}\E\norm{\nabla f(x_t)}^2
    \le\;&
    \frac{2(f(x_0)-f^\star)}{\eta T}
    +
    2\left(\frac{\sigma^2}{H}+\xi^2\right)
    \nonumber\\
    &+
    \frac{2}{T}\sum_{t=0}^{T-1}\E\norm{e_t}^2
    +
    L\eta M_{\tilde\Delta}^2.
    \label{eq:main_pre}
\end{align}
Define the average HRAC error floor
\[
    \mathcal{E}_T
    :=
    \frac{6}{T}\sum_{t=0}^{T-1}\E\!\left[
        2\sum_{i\in\HH} w_i^t \norm{\hat{\Delta}_i^t-\Delta_i^t}^2
        +2\sum_{i\in\HH} w_i^t(\norm{r_i^t}-\tau_i^t)_+^2
        +\varepsilon_{w,t}^2 \max_{i\in\HH}\norm{\Delta_i^t}^2
        +\beta_t \sum_{j\in\BB} w_j^t \norm{\tilde{\Delta}_j^t}^2
    \right].
\]
Using Lemma~\ref{lem:agg_bound}, one obtains
\begin{align}
    \frac{1}{T}\sum_{t=0}^{T-1}\E\norm{\nabla f(x_t)}^2
    \le\;&
    \frac{2(f(x_0)-f^\star)}{\eta T}
    +
    2\left(\frac{\sigma^2}{H}+\xi^2\right)
    +
    L\eta M_{\tilde\Delta}^2
    +
    \mathcal{E}_T.
    \label{eq:main_full}
\end{align}
Therefore, under the stated effective-adversary conditions, HRAC converges
sublinearly to a neighbourhood of a
first-order stationary point.
The neighbourhood radius is determined by the stochastic variance,
client heterogeneity, clipping bias, honest-weight distortion, and the
remaining Byzantine weight mass.
\end{theorem}

\begin{proof}
Sum the one-step descent inequality~\eqref{eq:descent} from \(t=0\) to
\(T-1\):
\begin{align*}
    \E[f(x_T)]
    \le\;&
    f(x_0)
    -
    \frac{\eta}{2}\sum_{t=0}^{T-1}\E\norm{\nabla f(x_t)}^2
    \\
    &+
    \eta\sum_{t=0}^{T-1}\E\norm{e_t}^2
    +
    \eta T\left(\frac{\sigma^2}{H}+\xi^2\right)
    +
    \frac{L\eta^2}{2}\sum_{t=0}^{T-1}\E\norm{g_t}^2.
\end{align*}
Since \(f(x_T)\ge f^\star\), rearranging gives
\[
    \frac{\eta}{2}\sum_{t=0}^{T-1}\E\norm{\nabla f(x_t)}^2
    \le
    f(x_0)-f^\star
    +
    \eta\sum_{t=0}^{T-1}\E\norm{e_t}^2
    +
    \eta T\left(\frac{\sigma^2}{H}+\xi^2\right)
    +
    \frac{L\eta^2}{2}\sum_{t=0}^{T-1}\E\norm{g_t}^2.
\]
Dividing both sides by \(\frac{\eta T}{2}\) yields
\[
    \frac{1}{T}\sum_{t=0}^{T-1}\E\norm{\nabla f(x_t)}^2
    \le
    \frac{2(f(x_0)-f^\star)}{\eta T}
    +
    \frac{2}{T}\sum_{t=0}^{T-1}\E\norm{e_t}^2
    +
    2\left(\frac{\sigma^2}{H}+\xi^2\right)
    +
    \frac{L\eta}{T}\sum_{t=0}^{T-1}\E\norm{g_t}^2.
\]
By Lemma~\ref{lem:g_second_derived}, \(\E\norm{g_t}^2 \le M_{\tilde\Delta}^2\).
This proves~\eqref{eq:main_pre}.
Substituting Lemma~\ref{lem:agg_bound} into \eqref{eq:main_pre}
yields~\eqref{eq:main_full}.
\end{proof}

\begin{proposition}[Tightened HRAC floor under \(\beta_t\)-linked distortion]
\label{prop:tight_floor_beta2}
Let $G_{\max} = \max_{i\in\HH}\norm{\Delta_i^t}$. Assume there exist constants \(C_w,C_m\ge 0\) such that for all large
\(t\):
\begin{align}
    \varepsilon_{w,t} &\le C_w\,\beta_t,
    \label{eq:tight_epsw_beta}\\
    \sum_{j\in\BB} w_j^t \norm{\tilde{\Delta}_j^t}^2
    &\le C_m\,\beta_t\,G_{\max}^2.
    \label{eq:tight_byz_energy}
\end{align}
Then the last two terms in \eqref{eq:agg_error_final} satisfy
\begin{equation}
    3\varepsilon_{w,t}^2 \max_{i\in\HH}\norm{\Delta_i^t}^2
    +
    3\beta_t \sum_{j\in\BB} w_j^t \norm{\tilde{\Delta}_j^t}^2
    \le
    3(C_w^2+C_m)\,\beta_t^2\,G_{\max}^2.
    \label{eq:tight_beta2_floor}
\end{equation}
Hence the corresponding component of \(\mathcal E_T\) decays at
\(O(\beta_t^2)\) when \(\beta_t\to 0\).
\end{proposition}

\begin{proof}
The first term is bounded by
\(
3(C_w\beta_t)^2 G_{\max}^2
\)
using \eqref{eq:tight_epsw_beta}.
The second term is bounded by
\(
3\beta_t(C_m\beta_t G_{\max}^2)
\)
using \eqref{eq:tight_byz_energy}. Summing gives
\eqref{eq:tight_beta2_floor}.
\end{proof}

\begin{corollary}[Refined asymptotic floor under \(\beta_t\)-linked distortion]
\label{cor:refined_beta2_floor}
Under Theorem~\ref{thm:main} and Proposition~\ref{prop:tight_floor_beta2},
\eqref{eq:main_full} admits the refinement
\[
    \mathcal{E}_T
    =
    \frac{6}{T}\sum_{t=0}^{T-1}\E\!\left[
        2\sum_{i\in\HH} w_i^t \norm{\hat{\Delta}_i^t-\Delta_i^t}^2
        +2\sum_{i\in\HH} w_i^t(\norm{r_i^t}-\tau_i^t)_+^2
        +\mathcal{O}(\beta_t^2)G_{\max}^2
    \right].
\]
In particular, the contribution of weight distortion and Byzantine mass in
\(\mathcal{E}_T\) decays quadratically in \(\beta_t\).
\end{corollary}

\begin{corollary}[Asymptotic stationarity under vanishing HRAC error]
\label{cor:vanishing}
Under the assumptions of Theorem~\ref{thm:main}, if additionally
\begin{align}
    \frac{1}{T}\sum_{t=0}^{T-1}\sum_{i\in\HH}
    \E\!\left[w_i^t \norm{\hat{\Delta}_i^t-\Delta_i^t}^2\right]
    &\to 0,
    \label{eq:v1}\\
    \frac{1}{T}\sum_{t=0}^{T-1}\sum_{i\in\HH}
    \E\!\left[w_i^t(\norm{r_i^t}-\tau_i^t)_+^2\right]
    &\to 0,
    \label{eq:v2}\\
    \frac{1}{T}\sum_{t=0}^{T-1}
    \E\!\left[\beta_t^2 G_{\max}^2\right]
    &\to 0,
    \label{eq:v3}
\end{align}
and the step size sequence satisfies \(\eta \to 0\) with \(T\eta \to \infty\),
then
\begin{equation}
    \frac{1}{T}\sum_{t=0}^{T-1}\E\norm{\nabla f(x_t)}^2 \to 0.
    \label{eq:stationarity}
\end{equation}
\end{corollary}

\begin{proof}
This follows immediately from~\eqref{eq:main_full}: all error-floor terms
vanish asymptotically, and the optimisation term
\((f(x_0)-f^\star)/(\eta T)\) also vanishes because \(T\eta \to \infty\).
\end{proof}
\begin{comment}
\section{Algorithmic Dynamics and Weight Control}

This section makes explicit use of the implemented recursion
\[
    \nu_i^{t+1} = \rho_\nu \nu_i^t + (1-\rho_\nu)d_i^t,
    \qquad 0<\rho_\nu<1.
\]

\begin{lemma}[Closed-form representation of \(\nu_i^t\)]
\label{lem:nu_closed_form}
For each client \(i\) and \(t\ge 1\),
\begin{equation}
    \nu_i^t
    =
    \rho_\nu^t \nu_i^0
    +
    (1-\rho_\nu)\sum_{k=0}^{t-1}\rho_\nu^k d_i^{t-1-k}.
    \label{eq:nu_closed_form}
\end{equation}
\end{lemma}

\begin{proof}
Unroll the linear recursion by induction.
\end{proof}

\paragraph{The statistic \(d_i^t\) (Implementation-aligned definition).}
The directional change statistic is computed directly from the consecutive
clipped residuals on the statistics path:
\begin{equation}
    d_i^t := \norm{\bar{r}_{i,\mathrm{norm}}^t - \bar{r}_{i,\mathrm{norm}}^{t-1}}.
    \label{eq:d_def}
\end{equation}
Here, \(\bar{r}_{i,\mathrm{norm}}^t\) denotes the clipped residual on the
statistics path. When no small-norm lifting is triggered, it coincides with the
ordinary clipped residual. When lifting is active, it is computed from the
lifted effective update used only by the statistics path. Thus, \(d_i^t\)
records the cross-round Euclidean displacement of the statistics-path clipped
residuals and serves as the direct input to the \(\nu\)-EMA recursion.

\begin{remark}[Small-Norm Lifting on the Statistics Path]
When a client's update norm is abnormally small relative to
the current round median norm, HRAC optionally rescales that update to the
median-norm level \emph{only on the statistics path}, in order to avoid
degenerate vanishing behaviour in the EMA statistics. This lifting affects the
computation of \(\mu_i^t\), \(\nu_i^t\), and the associated statistics-path
history variables, but does not directly alter the aggregation path or the
final per-round aggregated update.
\end{remark}

\begin{lemma}[EMA tracking of a non-vanishing \(d\)-band]
\label{lem:nu_tracking}
Let \(0<\rho_\nu<1\) and let
\[
    \nu_i^{t+1}
    =
    \rho_\nu \nu_i^t + (1-\rho_\nu)d_i^t.
\]
Then the following hold for every client \(i\):
\begin{enumerate}[label=(\roman*)]
    \item if \(d_i^t \to d_i^\star\), then \(\nu_i^t \to d_i^\star\);
    \item if there exist constants \(0\le \underline d_i \le \overline d_i < \infty\)
    and \(T_i<\infty\) such that
    \[
        \underline d_i \le d_i^t \le \overline d_i,
        \qquad \forall t\ge T_i,
    \]
    then
    \[
        \underline d_i
        \le
        \liminf_{t\to\infty}\nu_i^t
        \le
        \limsup_{t\to\infty}\nu_i^t
        \le
        \overline d_i.
    \]
\end{enumerate}
In particular, if \(d_i^t\) settles to a positive plateau or fluctuates inside
a positive bounded band, then \(\nu_i^t\) tracks the same positive late-stage
band rather than collapsing to \(0\).
\end{lemma}

\begin{proof}
For part (i), let \(e_i^t:=\nu_i^t-d_i^\star\). Then
\[
    e_i^{t+1}
    =
    \rho_\nu e_i^t + (1-\rho_\nu)(d_i^t-d_i^\star).
\]
Since \(d_i^t-d_i^\star\to 0\) and \(0<\rho_\nu<1\), the standard stable
linear recursion implies \(e_i^t\to 0\), hence \(\nu_i^t\to d_i^\star\).

For part (ii), for all \(t\ge T_i\),
\[
    \rho_\nu \nu_i^t + (1-\rho_\nu)\underline d_i
    \le
    \nu_i^{t+1}
    \le
    \rho_\nu \nu_i^t + (1-\rho_\nu)\overline d_i.
\]
Comparing with the scalar affine recursions
\[
    \underline v^{t+1}
    =
    \rho_\nu \underline v^t + (1-\rho_\nu)\underline d_i,
    \qquad
    \overline v^{t+1}
    =
    \rho_\nu \overline v^t + (1-\rho_\nu)\overline d_i,
\]
we obtain by induction that \(\nu_i^t\) is eventually trapped between two
sequences converging respectively to \(\underline d_i\) and \(\overline d_i\).
Taking \(\liminf\) and \(\limsup\) yields the claim.
\end{proof}

\begin{lemma}[Honest dual-path tube bound with explicit \(\eta\)-dependence]
\label{lem:honest_tube_eta}
Fix an honest client \(i\). Suppose:
\begin{enumerate}[label=(\roman*)]
    \item the honest local update decomposes as
    \(\Delta_i^t=\nabla f_i(x_t)+\zeta_i^t\) with
    \(\E[\zeta_i^t\mid\F_t]=0\),
    \(\E\!\left[\norm{\zeta_i^t}^2\mid\F_t\right]\le \sigma_i^2\);
    \item \(f_i\) is \(L_i\)-smooth:
    \(\norm{\nabla f_i(x_t)-\nabla f_i(x_{t-1})}\le L_i\norm{x_t-x_{t-1}}\);
    \item \(\sup_t \E\norm{g_t}^2 \le M_{\tilde\Delta}^2\), and \(x_{t+1}=x_t-\eta g_t\)
    with constant \(\eta\);
    \item the stats-path center uses EMA:
    \(b_{i,t+1}^{\mathrm{norm}}=\rho_b b_{i,t}^{\mathrm{norm}}+(1-\rho_b)q_{i,t}\),
    where \(\E\!\left[\norm{q_{i,t}-\Delta_i^t}^2\mid\F_t\right]\le \varepsilon_q^2\).
\end{enumerate}
Then there exists a finite constant \(C_i\) such that for all large \(t\),
\begin{equation}
    \E\!\left[\norm{\Delta_i^t-b_{i,t}^{\mathrm{norm}}}^2\right]
    \le
    \frac{C_i}{1-\rho_b}
    \Bigl(
        \sigma_i^2 + \eta^2 L_i^2 M_{\tilde\Delta}^2 + \varepsilon_q^2
    \Bigr).
    \label{eq:honest_tube_eta}
\end{equation}
In particular, as \(\eta\to 0\), the drift contribution is
\(O(\eta^2)\), so the tube radius is noise-dominated.
\end{lemma}

\begin{proof}
Define \(e_{i,t}:=\Delta_i^t-b_{i,t}^{\mathrm{norm}}\). By the EMA recursion,
\[
    e_{i,t+1}
    =
    \Delta_i^{t+1}
    -
    \rho_b b_{i,t}^{\mathrm{norm}}
    -
    (1-\rho_b)q_{i,t}
    =
    \rho_b e_{i,t}
    +
    (\Delta_i^{t+1}-\Delta_i^t)
    +
    (1-\rho_b)(\Delta_i^t-q_{i,t}).
\]
Set
\[
    w_{i,t}
    :=
    (\Delta_i^{t+1}-\Delta_i^t)
    +
    (1-\rho_b)(\Delta_i^t-q_{i,t}),
\]
so \(e_{i,t+1}=\rho_b e_{i,t}+w_{i,t}\).
For any \(\kappa>0\), Young's inequality gives
\[
    \norm{e_{i,t+1}}^2
    \le
    (1+\kappa)\rho_b^2\norm{e_{i,t}}^2
    +
    \left(1+\frac1\kappa\right)\norm{w_{i,t}}^2.
\]
Taking expectations,
\[
    \E\norm{e_{i,t+1}}^2
    \le
    a_\kappa \E\norm{e_{i,t}}^2
    +
    b_\kappa \E\norm{w_{i,t}}^2,
    \qquad
    a_\kappa:=(1+\kappa)\rho_b^2,\quad
    b_\kappa:=1+\frac1\kappa.
\]
Choose \(\kappa\in(0,\rho_b^{-2}-1)\), so \(a_\kappa<1\).
Now
\[
    \Delta_i^{t+1}-\Delta_i^t
    =
    \bigl(\nabla f_i(x_{t+1})-\nabla f_i(x_t)\bigr)
    +
    (\zeta_i^{t+1}-\zeta_i^t),
\]
so, by smoothness and \(x_{t+1}-x_t=-\eta g_t\),
\[
    \E\norm{\Delta_i^{t+1}-\Delta_i^t}^2
    \le
    2\eta^2L_i^2\E\norm{g_t}^2 + 4\sigma_i^2
    \le
    2\eta^2L_i^2 M_{\tilde\Delta}^2 + 4\sigma_i^2.
\]
Also, by \(\norm{u+v}^2\le 2\norm u^2+2\norm v^2\),
\[
    \E\norm{w_{i,t}}^2
    \le
    2\E\norm{\Delta_i^{t+1}-\Delta_i^t}^2
    +
    2(1-\rho_b)^2\varepsilon_q^2
    \le
    4\eta^2L_i^2M_{\tilde\Delta}^2 + 8\sigma_i^2 + 2(1-\rho_b)^2\varepsilon_q^2.
\]
Therefore
\[
    \E\norm{e_{i,t+1}}^2
    \le
    a_\kappa \E\norm{e_{i,t}}^2
    +
    b_\kappa
    \left(
        4\eta^2L_i^2M_{\tilde\Delta}^2 + 8\sigma_i^2 + 2(1-\rho_b)^2\varepsilon_q^2
    \right).
\]
Unrolling yields a geometric bound with factor \((1-a_\kappa)^{-1}\), hence
\eqref{eq:honest_tube_eta} with a finite constant \(C_i\).
\end{proof}

\begin{lemma}[EMA mismatch from input increments]
\label{lem:ema_increment_mismatch}
Let \(\{a_t\}_{t\ge 0}\subset\R^d\) and
\(
b_{t+1}=\rho b_t+(1-\rho)a_t
\)
with \(0<\rho<1\). Define \(e_t:=a_t-b_t\). Then
\begin{equation}
    e_{t+1}=\rho e_t + (a_{t+1}-a_t).
    \label{eq:ema_error_recursion}
\end{equation}
Consequently, for all \(t\ge 1\),
\begin{equation}
    \norm{e_t}
    \le
    \rho^t\norm{e_0}
    +
    \sum_{k=0}^{t-1}\rho^{\,t-1-k}\norm{a_{k+1}-a_k}.
    \label{eq:ema_error_unroll}
\end{equation}
In particular, if \(\norm{a_{t+1}-a_t}\to 0\), then \(\norm{e_t}\to 0\).
\end{lemma}

\begin{proof}
From
\(
e_{t+1}
=
a_{t+1}-b_{t+1}
=
a_{t+1}-\rho b_t-(1-\rho)a_t
\),
insert and subtract \(\rho a_t\):
\(
e_{t+1}
=
\rho(a_t-b_t)+(a_{t+1}-a_t)
=
\rho e_t+(a_{t+1}-a_t)
\),
which is \eqref{eq:ema_error_recursion}. Unrolling yields
\eqref{eq:ema_error_unroll}. If \(\norm{a_{t+1}-a_t}\to 0\), then the weighted
geometric average in \eqref{eq:ema_error_unroll} vanishes, and \(\rho^t\norm{e_0}\to 0\).
\end{proof}

\subsection{Adversarial Threat Model and Algorithmic Volatility Pressure}

\begin{lemma}[Expected honest \(d\)-upper bound from smooth tracking]
\label{lem:honest_d_upper}
Fix an honest client \(i\). Assume the hypotheses of
Lemma~\ref{lem:honest_tube_eta} hold and, in addition, the statistics-path
clipping map is non-expansive:
\[
    \norm{\bar r_{i,\mathrm{norm}}(u)-\bar r_{i,\mathrm{norm}}(v)}
    \le
    \norm{u-v},
    \qquad
    \forall u,v\in\R^d.
\]
Then there exists a finite constant \(r_{H,i}\) such that for all sufficiently
large \(t\),
\begin{equation}
    \E[d_i^t]
    \le
    r_{H,i},
    \label{eq:honest_d_upper}
\end{equation}
with
\[
    r_{H,i}
    =
    \mathcal O\!\left(
        \eta L_i M_{\tilde\Delta} + \sigma_i + \varepsilon_q
    \right).
\]
Consequently, with
\(
r_H:=\max_{i\in\HH} r_{H,i},
\)
\begin{equation}
    \max_{i\in\HH}\E[d_i^t] \le r_H
    \qquad \text{for all sufficiently large } t.
    \label{eq:honest_d_upper_global}
\end{equation}
\end{lemma}

\begin{proof}[Proof idea]
Write
\(
\bar r_{i,\mathrm{norm}}^t
=
\bar r_{i,\mathrm{norm}}(\Delta_i^t-b_{i,t}^{\mathrm{norm}})
\).
By non-expansiveness,
\[
    d_i^t
    \le
    \norm{
        (\Delta_i^t-b_{i,t}^{\mathrm{norm}})
        -
        (\Delta_i^{t-1}-b_{i,t-1}^{\mathrm{norm}})
    }.
\]
Hence
\[
    d_i^t
    \le
    \norm{\Delta_i^t-\Delta_i^{t-1}}
    +
    \norm{b_{i,t}^{\mathrm{norm}}-b_{i,t-1}^{\mathrm{norm}}}.
\]
The first term is controlled by smoothness and bounded stochastic variation:
\[
    \E\norm{\Delta_i^t-\Delta_i^{t-1}}
    =
    \mathcal O(\eta L_i M_{\tilde\Delta}+\sigma_i).
\]
For the second term, use the EMA update of \(b_{i,t}^{\mathrm{norm}}\):
\[
    b_{i,t}^{\mathrm{norm}}-b_{i,t-1}^{\mathrm{norm}}
    =
    (1-\rho_b)\bigl(q_{i,t-1}-b_{i,t-1}^{\mathrm{norm}}\bigr),
\]
so Lemma~\ref{lem:honest_tube_eta} yields
\[
    \E\norm{b_{i,t}^{\mathrm{norm}}-b_{i,t-1}^{\mathrm{norm}}}
    =
    \mathcal O(\eta L_i M_{\tilde\Delta}+\sigma_i+\varepsilon_q).
\]
Taking expectation in the bound on \(d_i^t\) and combining the two estimates
gives \eqref{eq:honest_d_upper}; taking the maximum over honest clients yields
\eqref{eq:honest_d_upper_global}. This lemma is therefore an expectation-level
smoothness estimate. Any pathwise eventual bound on honest \(d_i^t\) used later
must be imposed separately or derived from additional concentration arguments.
\end{proof}

Instead of assuming arbitrary Byzantine behavior or postulating an artificial
\(d\)-gap, we isolate an HRAC-specific mechanism: persistent adversarial
history mismatch leaves an observable trace on the statistics path and thereby
creates a route to higher late-stage Byzantine \(d\)- and \(\nu\)-bands. The
empirical target is relative elevation, not a literal honest-to-zero versus
Byzantine-nonzero dichotomy.

\begin{definition}[Persistently effective attack]
\label{def:effective_attack}
An adversarial strategy is \emph{persistently effective} on a late-stage time
window if it induces a non-vanishing Byzantine contribution to the aggregate
throughout that window. Formally, there exist \(T_\epsilon<\infty\),
\(\epsilon > 0\), and a window length \(W\ge 1\) such that
\[
    \norm{M_t}
    =
    \norm{\sum_{j\in\BB} w_j^t \tilde{\Delta}_j^t}
    \ge \epsilon,
    \qquad
    \forall t\in\{T_\epsilon,\dots,T_\epsilon+W-1\}.
\]
\end{definition}

\begin{condition}[Persistent mismatch with observable residual motion]
\label{cond:persistent_mismatch_motion}
Fix a Byzantine client \(j\in\BB\). Let \(r_{j,\mathrm{norm}}^t\) denote the
pre-clipped residual on the statistics path and
\(
\bar r_{j,\mathrm{norm}}^t
=
\clip(r_{j,\mathrm{norm}}^t,\tau_{j,\mathrm{norm}}^t)
\)
the corresponding clipped residual. We say that \(j\) exhibits
\emph{persistent mismatch with observable residual motion} if there exist
\(T_B<\infty\), \(\underline\delta_B>0\), and \(\lambda_B\in(0,1]\) such that
for all \(t\ge T_B\),
\begin{align}
    \norm{r_{j,\mathrm{norm}}^t-r_{j,\mathrm{norm}}^{t-1}}
    &\ge \underline\delta_B,
    \label{eq:mismatch_motion_raw_lb}\\
    \norm{\bar r_{j,\mathrm{norm}}^t-\bar r_{j,\mathrm{norm}}^{t-1}}
    &\ge
    \lambda_B
    \norm{r_{j,\mathrm{norm}}^t-r_{j,\mathrm{norm}}^{t-1}}.
    \label{eq:mismatch_motion_clip_lb}
\end{align}
\end{condition}

\begin{proposition}[Conditional route from persistent mismatch to a higher Byzantine \(d\)-band]
\label{prop:mismatch_to_d_band}
Assume Condition~\ref{cond:persistent_mismatch_motion} holds for a Byzantine
client \(j\in\BB\). Then, for all \(t\ge T_B\),
\begin{equation}
    d_j^t
    =
    \norm{\bar r_{j,\mathrm{norm}}^t-\bar r_{j,\mathrm{norm}}^{t-1}}
    \ge
    \lambda_B \underline\delta_B
    =: r_{B,j}^{\mathrm{mis}}.
    \label{eq:mismatch_to_d_lb}
\end{equation}
Consequently, if in addition the honest late-stage upper bound
\eqref{eq:honest_d_cond} holds and
\begin{equation}
    r_{B,j}^{\mathrm{mis}} > r_H,
    \label{eq:mismatch_margin}
\end{equation}
then client \(j\) satisfies
\begin{equation}
    d_j^t-\max_{i\in\HH} d_i^t
    \ge
    r_{B,j}^{\mathrm{mis}}-r_H
    >0,
    \qquad \forall t\ge \max\{T_B,T_H\}.
    \label{eq:mismatch_pointwise_d_gap}
\end{equation}
Thus persistent mismatch yields a formal sufficient route to a higher
Byzantine late-stage \(d\)-band, and hence to the pointwise \(\nu\)-route
developed below.
\end{proposition}

\begin{proof}
By \eqref{eq:mismatch_motion_clip_lb} and
\eqref{eq:mismatch_motion_raw_lb},
\[
    d_j^t
    =
    \norm{\bar r_{j,\mathrm{norm}}^t-\bar r_{j,\mathrm{norm}}^{t-1}}
    \ge
    \lambda_B
    \norm{r_{j,\mathrm{norm}}^t-r_{j,\mathrm{norm}}^{t-1}}
    \ge
    \lambda_B\underline\delta_B,
\]
which proves \eqref{eq:mismatch_to_d_lb}. If also
\eqref{eq:honest_d_cond} holds, then for all \(t\ge \max\{T_B,T_H\}\),
\[
    d_j^t-\max_{i\in\HH}d_i^t
    \ge
    r_{B,j}^{\mathrm{mis}}-r_H.
\]
Under \eqref{eq:mismatch_margin}, this is strictly positive, giving
\eqref{eq:mismatch_pointwise_d_gap}.
\end{proof}

\begin{remark}[Interpretation of Condition~\ref{cond:persistent_mismatch_motion}]
\label{rem:mismatch_motion_interp}
Condition~\ref{cond:persistent_mismatch_motion} isolates the weakest step in a
form that is strong enough for the later proofs but still faithful to the
implementation. The first inequality requires that the attack keep generating
nontrivial cross-round motion on the statistics-path residual before clipping.
The second requires that clipping not completely erase that motion. Together
they turn qualitative persistent mismatch into a quantitative lower bound on
\(d_j^t\).
\end{remark}

\begin{remark}[Mechanism-level reading]
\label{rem:volatility_pressure}
The code-aligned interpretation is that a Byzantine update stream which stays
persistently mismatched with HRAC's local history state should satisfy
Condition~\ref{cond:persistent_mismatch_motion} on late-stage windows whenever
the mismatch keeps reappearing and is not annihilated by clipping. In that
case, Proposition~\ref{prop:mismatch_to_d_band} formalizes the route from
persistent mismatch to a higher Byzantine \(d\)-band.
\end{remark}

\begin{remark}[Equivalent informal picture]
Under the same mechanism, the statistics-path clipped residual increments
\[
    d_j^t
    =
    \norm{\bar r_{j,\mathrm{norm}}^t-\bar r_{j,\mathrm{norm}}^{t-1}}
\]
remain elevated relative to the honest late-stage tracking band over long horizons, so
that Byzantine clients occupy a higher late-stage \(d\)-band rather than the
same lower band as honest clients. The corresponding EMA statistic \(\nu_j^t\)
then tracks this enlarged late-stage band.
\end{remark}

\begin{remark}[Label-flipping as a concrete code-aligned example]
\label{rem:label_flipping_example}
The label-flipping attack used in our codebase provides a direct illustration
of Remark~\ref{rem:volatility_pressure} and a plausible route to
Condition~\ref{cond:persistent_mismatch_motion}. In that attack, a Byzantine
client does not merely add independent noise to an otherwise honest local
objective; instead, it trains on systematically corrupted labels. As a result,
its local update direction is generated by a different effective objective than
that of the honest clients. From the viewpoint of HRAC, this means that the
attacker's effective update on the statistics path is repeatedly pulled away
from the history state that was built from previous clipped messages.

Consequently, the attacker is less likely to remain inside the same
comparatively smooth and lower late-stage tracking band as an honest client.
Even after global clipping and
residual clipping bound the per-round magnitude, the cross-round displacement
\[
    d_j^t
    =
    \norm{\bar r_{j,\mathrm{norm}}^t-\bar r_{j,\mathrm{norm}}^{t-1}}
\]
tends to stay larger because the attack keeps reintroducing history mismatch.
The EMA recursion
\[
    \nu_j^{t+1}
    =
    \rho_\nu \nu_j^t + (1-\rho_\nu)d_j^t
\]
then accumulates this larger late-stage \(d_j^t\)-band into a larger
late-stage \(\nu_j^t\)-band. This is the code-aligned reason that, in the
label-flipping runs, Byzantine clients need not drive \(\nu_j^t\) to an
extreme value at every round, but they do tend to occupy a visibly higher
\(\nu\)-band than honest clients over long horizons.
\end{remark}

\begin{proposition}[Auxiliary zero-absorption criterion for persistently effective attacks]
\label{prop:attack_nonabsorb}
Assume a Byzantine client \(j\) participates in a persistently effective attack
with magnitude \(\epsilon>0\), in the sense of Definition~\ref{def:effective_attack},
and that its statistics-path effective update remains persistently mismatched
with the HRAC history-tracking state. Assume moreover:
\begin{enumerate}[label=(\roman*)]
    \item the per-round malicious innovation is bounded by HRAC clipping and the
    global cap, so the attack can persist only through its residual mismatch
    with the tracked history state;
    \item if \(d_j^t\to 0\), then the statistics-path clipped residual state
    \(\bar r_{j,\mathrm{norm}}^t\) becomes asymptotically stationary relative to
    that tracked state, so the attack-induced mismatch is asymptotically
    eliminated;
    \item once this mismatch is eliminated, the induced injected innovation
    against the evolving model decays below the effectiveness threshold
    \(\epsilon\).
\end{enumerate}
Then there exists a positive function \(r_B=r_B(\epsilon,c,\rho_b)\) such that
\begin{equation}
    \liminf_{t\to\infty} d_j^t \ge r_B(\epsilon,c,\rho_b)>0.
    \label{eq:byz_d_lower}
\end{equation}
\end{proposition}

\begin{proof}[Proof idea]
Suppose, to the contrary, that \(\liminf_{t\to\infty} d_j^t = 0\). Then there
exists a subsequence along which the clipped residual increments become
arbitrarily small. By assumption (ii), this drives the statistics-path clipped
residual state \(\bar r_{j,\mathrm{norm}}^t\) into an asymptotically stationary
regime relative to the tracked history state. The attack-induced mismatch is
then progressively absorbed by the HRAC history dynamics, leaving only a
vanishing residual innovation after clipping. Under assumptions (i) and (iii),
this prevents the attacker from maintaining a persistent aggregate bias of size
at least \(\epsilon\), contradicting
Definition~\ref{def:effective_attack}. Therefore \(d_j^t\) cannot approach
zero, and the positive lower bound \eqref{eq:byz_d_lower} follows.
\end{proof}

\begin{remark}[Interpretation of Proposition~\ref{prop:attack_nonabsorb}]
Proposition~\ref{prop:attack_nonabsorb} formalizes the code-aligned intuition
that persistent adversarial mismatch with the tracked history state is visible
through the \(d_j^t\)-statistics. The proposition is not claiming that an
attacker explicitly chooses to maximize \(d_j^t\). Rather, it states that for
the attack classes considered here, elevated \(d_j^t\) is a dynamical
consequence of staying adversarially misaligned with the comparatively smooth
honest late-stage tracking band.
\end{remark}

\begin{remark}[Role of Proposition~\ref{prop:attack_nonabsorb} in the present manuscript]
Proposition~\ref{prop:attack_nonabsorb} should be read as an auxiliary
sufficient route, not as the primary empirical claim. The experiments used in
this note do not suggest that honest \(d_i^t\) values converge to \(0\); on
the contrary, both honest and Byzantine clients remain in positive late-stage
bands. The practically relevant distinction is that Byzantine clients tend to
occupy a \emph{higher} \(d\)- and therefore \(\nu\)-band, which is exactly the
relative-elevation effect used by the weighting rule.
\end{remark}

\begin{remark}[Scope of the adaptive-statistics layer]
\label{rem:adaptive_scope}
The adaptive-statistics argument has three different proof strengths. First,
the recursion-level statements, such as Lemma~\ref{lem:nu_tracking},
Theorem~\ref{thm:mean_d_to_mean_nu}, and
Theorem~\ref{thm:mean_nu_to_penalty}, are exact consequences of the implemented
EMA and weighting formulas. Second, results such as
Theorem~\ref{thm:nu_separation_main} are conditional upgrades that require
additional late-stage pathwise control. Third,
Condition~\ref{cond:persistent_mismatch_motion},
Proposition~\ref{prop:mismatch_to_d_band}, and
Proposition~\ref{prop:attack_nonabsorb} are intentionally mechanism-level:
they formalize a code-aligned route from persistent history mismatch to a
higher Byzantine \(d\)-band, but they are not meant to replace a full
attack-specific stochastic analysis. This is why the manuscript uses the
group-mean route as the main storyline and treats the pointwise route as a
stronger sufficient refinement.
\end{remark}

\begin{condition}[Optional late-stage honest \(d\)-upper bound]
\label{cond:honest_d_upper}
There exist \(T_H<\infty\) and \(r_H\ge 0\) such that
\begin{equation}
    \max_{i\in\HH} d_i^t \le r_H,
    \qquad \forall t\ge T_H.
    \label{eq:honest_d_cond}
\end{equation}
\end{condition}

\begin{remark}[How Lemma~\ref{lem:honest_d_upper} relates to Condition~\ref{cond:honest_d_upper}]
Lemma~\ref{lem:honest_d_upper} provides an expectation-level estimate for the
honest \(d\)-statistics. Condition~\ref{cond:honest_d_upper} is the stronger
pathwise late-stage form needed by the pointwise \(\nu\)-separation argument.
It may be justified by additional concentration assumptions, or treated as a
code-aligned sufficient condition when one works directly with the realized
training trajectory.
\end{remark}

\begin{remark}[Status of the pointwise route]
\label{rem:pointwise_route_status}
The next theorem is kept as a stronger sufficient route, but it is not the
main storyline of this manuscript. Its role is to show that, under additional
pathwise control, one can upgrade relative late-stage band elevation to a full
pointwise Byzantine--honest \(\nu\)-gap. The experimentally primary route
remains the weaker group-mean mechanism developed below.
\end{remark}

\begin{theorem}[Stronger pointwise route from honest \(d\)-control and attack effectiveness]
\label{thm:nu_separation_main}
Assume Condition~\ref{cond:honest_d_upper} holds and that every Byzantine
client \(j\in\BB\) satisfies the conclusion of
Proposition~\ref{prop:attack_nonabsorb} with the same positive lower bound. Let
\[
    \qquad
    r_B := r_B(\epsilon,c,\rho_b).
\]
If
\begin{equation}
    r_B > r_H,
\label{eq:strict_d_margin}
\end{equation}
then there exist \(\underline\Delta_d:=\frac12(r_B-r_H)>0\) and \(T_0<\infty\)
such that for all
sufficiently large \(t\),
\begin{equation}
    \min_{j\in\BB} d_j^t - \max_{i\in\HH} d_i^t \ge \underline\Delta_d.
    \label{eq:eventual_d_gap}
\end{equation}
Consequently, for every pair \((j,i)\in\BB\times\HH\),
\begin{equation}
    \liminf_{t\to\infty}(\nu_j^t-\nu_i^t)\ge \underline\Delta_d,
    \label{eq:eventual_nu_gap_pair}
\end{equation}
and therefore there exists \(T_1<\infty\) such that Condition~\ref{cond:nu_sep_emp}
holds with some \(0<\gamma_\nu<\underline\Delta_d\) and finite spreads
\(\omega_H,\omega_B\).
\end{theorem}

\begin{proof}
By Condition~\ref{cond:honest_d_upper},
\[
    \max_{i\in\HH} d_i^t \le r_H
    \qquad \forall t\ge T_H.
\]
By the uniform Byzantine lower-bound hypothesis,
\[
    \liminf_{t\to\infty}\min_{j\in\BB} d_j^t \ge r_B.
\]
Since \(r_B-r_H>0\), by the definition of \(\liminf\) there exists
\(T_B<\infty\) such that
\[
    \min_{j\in\BB} d_j^t \ge \frac{r_B+r_H}{2},
    \qquad \forall t\ge T_B.
\]
Therefore, for all \(t\ge T_0:=\max\{T_H,T_B\}\),
\[
    \min_{j\in\BB} d_j^t - \max_{i\in\HH} d_i^t
    \ge
    \frac{r_B+r_H}{2} - r_H
    =
    \frac{r_B-r_H}{2}
    =
    \underline\Delta_d.
\]
This proves
\eqref{eq:eventual_d_gap}.

Now fix any \((j,i)\in\BB\times\HH\). By Lemma~\ref{lem:nu_gap_recursion},
\[
    z_{ji}^{t+1}
    =
    \rho_\nu z_{ji}^t + (1-\rho_\nu)(d_j^t-d_i^t)
    \ge
    \rho_\nu z_{ji}^t + (1-\rho_\nu)\underline\Delta_d
\]
for all \(t\ge T_0\). Unrolling gives
\[
    z_{ji}^{t}
    \ge
    \rho_\nu^{\,t-T_0} z_{ji}^{T_0}
    +
    (1-\rho_\nu^{\,t-T_0})\underline\Delta_d.
\]
Taking \(\liminf\) yields \eqref{eq:eventual_nu_gap_pair}. Since the honest and
Byzantine \(\nu\)-trajectories are bounded by Lemma~\ref{lem:nu_tracking}
applied to bounded \(d\)-bands, their within-group spreads remain finite on
late-stage windows. Hence one may choose \(T_1\) large enough so that the
pointwise gap remains above some \(0<\gamma_\nu<\underline\Delta_d\), which yields
Condition~\ref{cond:nu_sep_emp}.
\end{proof}

\begin{condition}[Late-stage \(\nu\)-band separation]
\label{cond:nu_sep_emp}
We say that late-stage \(\nu\)-band separation holds if there exist \(T_1<\infty\),
\(\gamma_\nu>0\), \(\omega_H\ge 0\), and \(\omega_B\ge 0\) such that for all
\(t\ge T_1\),
\begin{align}
    \min_{j\in\BB}\nu_j^t - \max_{i\in\HH}\nu_i^t &\ge \gamma_\nu,
    \label{eq:nu_gap_ass}\\
    \max_{i\in\HH}\nu_i^t - \min_{i\in\HH}\nu_i^t &\le \omega_H,
    \label{eq:nu_h_spread}\\
    \max_{j\in\BB}\nu_j^t - \min_{j\in\BB}\nu_j^t &\le \omega_B.
    \label{eq:nu_b_spread}
\end{align}
\end{condition}

\begin{definition}[Group means and radii for \(\nu\)]
\label{def:group_mean_radius}
For each round \(t\), define
\[
    \bar\nu_H^t := \frac1H\sum_{i\in\HH}\nu_i^t,
    \qquad
    \bar\nu_B^t := \frac1B\sum_{j\in\BB}\nu_j^t,
\]
and the corresponding within-group radii
\[
    r_H^t := \max_{i\in\HH}\lvert \nu_i^t-\bar\nu_H^t\rvert,
    \qquad
    r_B^t := \max_{j\in\BB}\lvert \nu_j^t-\bar\nu_B^t\rvert.
\]
\end{definition}

\begin{theorem}[Late-stage group-mean \(d\)-band elevation implies group-mean \(\nu\)-band elevation]
\label{thm:mean_d_to_mean_nu}
For each round \(t\), define the group means
\[
    \bar d_H^t := \frac1H\sum_{i\in\HH} d_i^t,
    \qquad
    \bar d_B^t := \frac1B\sum_{j\in\BB} d_j^t.
\]
Assume there exist \(T_d<\infty\) and \(\bar\gamma_d>0\) such that for all
\(t\ge T_d\),
\begin{equation}
    \bar d_B^t-\bar d_H^t \ge \bar\gamma_d.
    \label{eq:mean_d_sep}
\end{equation}
Then the group-mean \(\nu\)-gap satisfies the recursion
\begin{equation}
    \bar\nu_B^{t+1}-\bar\nu_H^{t+1}
    =
    \rho_\nu(\bar\nu_B^t-\bar\nu_H^t)
    +
    (1-\rho_\nu)(\bar d_B^t-\bar d_H^t),
    \label{eq:mean_nu_gap_recursion}
\end{equation}
and therefore, for all \(t\ge T_d\),
\begin{equation}
    \bar\nu_B^{t}-\bar\nu_H^{t}
    \ge
    \rho_\nu^{\,t-T_d}(\bar\nu_B^{T_d}-\bar\nu_H^{T_d})
    +
    (1-\rho_\nu^{\,t-T_d})\bar\gamma_d.
    \label{eq:mean_nu_gap_unrolled}
\end{equation}
In particular,
\begin{equation}
    \liminf_{t\to\infty}(\bar\nu_B^t-\bar\nu_H^t)\ge \bar\gamma_d>0.
    \label{eq:mean_nu_gap_from_mean_d}
\end{equation}
Hence Condition~\ref{cond:mean_nu_sep} holds with any
\(0<\bar\gamma_\nu<\bar\gamma_d\) after a finite transient.
\end{theorem}

\begin{proof}
Average the recursion
\(
\nu_i^{t+1}=\rho_\nu\nu_i^t+(1-\rho_\nu)d_i^t
\)
over \(i\in\HH\) and over \(j\in\BB\), obtaining
\[
    \bar\nu_H^{t+1}
    =
    \rho_\nu\bar\nu_H^t+(1-\rho_\nu)\bar d_H^t,
    \qquad
    \bar\nu_B^{t+1}
    =
    \rho_\nu\bar\nu_B^t+(1-\rho_\nu)\bar d_B^t.
\]
Subtracting the two identities gives \eqref{eq:mean_nu_gap_recursion}. Under
\eqref{eq:mean_d_sep}, for all \(t\ge T_d\),
\[
    \bar\nu_B^{t+1}-\bar\nu_H^{t+1}
    \ge
    \rho_\nu(\bar\nu_B^t-\bar\nu_H^t)+(1-\rho_\nu)\bar\gamma_d.
\]
Unrolling this affine recursion yields \eqref{eq:mean_nu_gap_unrolled}. Taking
\(\liminf\) proves \eqref{eq:mean_nu_gap_from_mean_d}. The final claim follows
because the transient term
\(\rho_\nu^{\,t-T_d}(\bar\nu_B^{T_d}-\bar\nu_H^{T_d})\) vanishes.
\end{proof}

\begin{remark}[Why Theorem~\ref{thm:mean_d_to_mean_nu} matches the experiments]
\label{rem:mean_d_to_mean_nu_emp}
Theorem~\ref{thm:mean_d_to_mean_nu} captures the empirically most faithful
route in the present manuscript. The logs and best-fit curves suggest that both
honest and Byzantine clients remain in positive late-stage \(d\)-bands, while
the Byzantine group occupies a higher group-mean \(d\)-band. The theorem then
shows that this relative \(d\)-band elevation is automatically accumulated by
the EMA recursion into a higher group-mean \(\nu\)-band, which is the quantity that
directly drives weighting.
\end{remark}

\begin{condition}[Late-stage group-mean \(\nu\)-band elevation]
\label{cond:mean_nu_sep}
We say that late-stage group-mean \(\nu\)-band elevation holds if there exist
\(T_{\mathrm{mean}}<\infty\) and \(\bar\gamma_\nu>0\) such that for all
\(t\ge T_{\mathrm{mean}}\),
\begin{equation}
    \bar\nu_B^t-\bar\nu_H^t \ge \bar\gamma_\nu > 0.
    \label{eq:mean_nu_sep}
\end{equation}
\end{condition}

\begin{theorem}[Late-stage Byzantine \(\nu\)-band elevation implies stronger penalty pressure]
\label{thm:mean_nu_to_penalty}
Assume Condition~\ref{cond:mean_nu_sep} and define the idealized excess scores
\(
e_k^t := (\nu_k^t-\bar\nu_t)_+
\)
with
\(
\bar\nu_t := \frac1N\sum_{\ell=1}^N \nu_\ell^t.
\)
Then for every \(t\ge T_{\mathrm{mean}}\),
\begin{equation}
    \frac1B\sum_{j\in\BB} e_j^t
    \ge
    \frac{H}{N}\,\bar\gamma_\nu.
    \label{eq:mean_excess_byz_lb}
\end{equation}
Consequently, at every such round there exists at least one Byzantine client
\(j_t\in\BB\) satisfying
\begin{equation}
    e_{j_t}^t \ge \frac{H}{N}\,\bar\gamma_\nu,
    \label{eq:exists_byz_penalized}
\end{equation}
and therefore at least one Byzantine raw score obeys
\begin{equation}
    \hat w_{j_t}^t
    =
    \frac{1}{1+\alpha e_{j_t}^t}
    \le
    \frac{1}{1+\alpha \frac{H}{N}\bar\gamma_\nu}.
    \label{eq:exists_byz_raw_weight_upper}
\end{equation}
\end{theorem}

\begin{proof}
Fix \(t\ge T_{\mathrm{mean}}\). Since
\[
    \bar\nu_t
    =
    \frac{H\bar\nu_H^t+B\bar\nu_B^t}{N}
    =
    \bar\nu_B^t-\frac{H}{N}(\bar\nu_B^t-\bar\nu_H^t),
\]
Condition~\ref{cond:mean_nu_sep} gives
\[
    \bar\nu_B^t-\bar\nu_t
    =
    \frac{H}{N}(\bar\nu_B^t-\bar\nu_H^t)
    \ge
    \frac{H}{N}\bar\gamma_\nu.
\]
Now use Jensen's inequality for the convex map \(x\mapsto x_+\):
\[
    \frac1B\sum_{j\in\BB} e_j^t
    =
    \frac1B\sum_{j\in\BB}(\nu_j^t-\bar\nu_t)_+
    \ge
    \left(
        \frac1B\sum_{j\in\BB}(\nu_j^t-\bar\nu_t)
    \right)_+
    =
    (\bar\nu_B^t-\bar\nu_t)_+.
\]
Combining the two displays yields \eqref{eq:mean_excess_byz_lb}. Hence the
average Byzantine excess is at least \(\frac{H}{N}\bar\gamma_\nu\), so at least
one Byzantine client satisfies \eqref{eq:exists_byz_penalized}. The raw-weight
bound \eqref{eq:exists_byz_raw_weight_upper} follows from monotonicity of
\(x\mapsto 1/(1+\alpha x)\).
\end{proof}

\begin{proposition}[Pointwise \(\nu\)-separation from mean separation and radii]
\label{prop:mean_radius_to_pointwise}
Suppose there exist \(T_1<\infty\), \(\bar\gamma_\nu>0\), \(r_H\ge 0\), and
\(r_B\ge 0\) such that for all \(t\ge T_1\),
\begin{equation}
    \bar\nu_B^t-\bar\nu_H^t \ge \bar\gamma_\nu,
    \qquad
    r_H^t \le r_H,
    \qquad
    r_B^t \le r_B.
    \label{eq:mean_radius_cond}
\end{equation}
If
\begin{equation}
    \bar\gamma_\nu > r_H+r_B,
    \label{eq:mean_radius_margin}
\end{equation}
then Condition~\ref{cond:nu_sep_emp} holds with
\begin{equation}
    \gamma_\nu = \bar\gamma_\nu-r_H-r_B,
    \qquad
    \omega_H = 2r_H,
    \qquad
    \omega_B = 2r_B.
    \label{eq:mean_radius_to_band}
\end{equation}
\end{proposition}

\begin{proof}
For \(t\ge T_1\),
\[
    \min_{j\in\BB}\nu_j^t
    \ge
    \bar\nu_B^t-r_B,
    \qquad
    \max_{i\in\HH}\nu_i^t
    \le
    \bar\nu_H^t+r_H.
\]
Hence
\[
    \min_{j\in\BB}\nu_j^t-\max_{i\in\HH}\nu_i^t
    \ge
    (\bar\nu_B^t-\bar\nu_H^t)-r_H-r_B
    \ge
    \bar\gamma_\nu-r_H-r_B
    =
    \gamma_\nu.
\]
Also,
\[
    \max_{i\in\HH}\nu_i^t-\min_{i\in\HH}\nu_i^t\le 2r_H,
    \qquad
    \max_{j\in\BB}\nu_j^t-\min_{j\in\BB}\nu_j^t\le 2r_B.
\]
Thus \eqref{eq:nu_gap_ass}--\eqref{eq:nu_b_spread} follow.
\end{proof}

\begin{remark}[Implementation Alignment]
Condition~\ref{cond:nu_sep_emp} matches the actual role played by the
implementation: HRAC does not require a pointwise eventual \(d\)-gap at every
round. What matters operationally is late-stage separation in the EMA
statistic \(\nu\), since \(\nu\) is the quantity that directly enters the
weighting rule. This is also the quantity exposed by the logs.
\end{remark}

\paragraph{Algorithmic Intuition: The Decoupled Penalty Mechanism.}
The implementation induces a structural trade-off for adversarial clients.
After an initial warm-up stage, HRAC assigns
unnormalized aggregation scores according to
\begin{equation}
    \hat{w}_i^t = \frac{1}{1+\alpha \, \xi_i^t},
    \qquad
    \xi_i^t =
    \begin{cases}
        \max(\nu_i^t-\overline{\nu}^t, 0), & \nu_i^t > 0.9,\\
        0, & \text{otherwise},
    \end{cases}
    \label{eq:exact_relative_penalty}
\end{equation}
followed by normalization
\(
w_i^t = \hat{w}_i^t / \sum_k \hat{w}_k^t
\),
possibly followed by iterative max-weight capping in implementation.

This decouples single-round mitigation from cross-round tracking:
\begin{enumerate}
    \item \textbf{Single-Round Mitigation:} The history center \(b_i^t\) and
    residual cap \(\tau_i^t\) explicitly bound the immediate malicious
    injection magnitude.
    \item \textbf{Cross-Round Penalty:} The \(\nu_i^t\) statistic evaluates
    behavioral stability through
    \[
        d_i^t
        =
        \norm{\bar r_{i,\mathrm{norm}}^t-\bar r_{i,\mathrm{norm}}^{t-1}}.
    \]
    Small \(d_i^t\) slows \(\nu\)-accumulation but also limits temporal
    adaptability of the injected residuals; large \(d_i^t\) improves
    adaptability but raises \(\nu_i^t\) relative to the population mean.
    Once \(\nu_i^t>0.9\), this increase is converted into down-weighting by
    \eqref{eq:exact_relative_penalty}.
\end{enumerate}

\begin{remark}[Positive-band interpretation]
The intended interpretation of \eqref{eq:exact_relative_penalty} is relative,
not absolute. Honest clients need not have \(\nu_i^t\approx 0\), and Byzantine
clients need not be separated by a zero-versus-nonzero dichotomy. The mechanism
only requires that Byzantine clients tend to exhibit a larger late-stage
\(\nu\)-band, so that their excess scores \((\nu_i^t-\bar\nu_t)_+\) are
activated more often or more strongly than those of honest clients.
\end{remark}

\subsection{Excess Gap and Ideal Weight Control}

\begin{lemma}[Pairwise \(\nu\)-gap recursion]
\label{lem:nu_gap_recursion}
For any \(j\in\BB\), \(i\in\HH\), define
\(z_{ji}^t:=\nu_j^t-\nu_i^t\). Then
\begin{equation}
    z_{ji}^{t+1}
    =
    \rho_\nu z_{ji}^t + (1-\rho_\nu)(d_j^t-d_i^t).
    \label{eq:z_recursion}
\end{equation}
\end{lemma}

\begin{proof}
Subtract the two \(\nu\)-updates:
\[
    \nu_j^{t+1}-\nu_i^{t+1}
    =
    \rho_\nu(\nu_j^t-\nu_i^t)+(1-\rho_\nu)(d_j^t-d_i^t).
\]
\end{proof}

\begin{lemma}[Exact \(m\)-step \(\nu\)-gap identity]
\label{lem:m_step_nu_gap}
For any \(m\ge 1\) and any pair \((j,i)\in\BB\times\HH\),
\begin{equation}
    z_{ji}^{t+m}
    =
    \rho_\nu^m z_{ji}^t
    +
    (1-\rho_\nu)\sum_{\ell=0}^{m-1}\rho_\nu^{m-1-\ell}
    \bigl(d_j^{t+\ell}-d_i^{t+\ell}\bigr).
    \label{eq:m_step_z_exact}
\end{equation}
\end{lemma}

\begin{proof}
Iterate \eqref{eq:z_recursion} \(m\) times and collect the resulting geometric
weights.
\end{proof}

\begin{definition}[Windowed weighted \(d\)-gap certificate]
\label{def:window_gap}
For \(m\ge 1\), define the \(m\)-step weighted \(d\)-gap between Byzantine
client \(j\) and honest client \(i\) by
\begin{equation}
    \mathcal G_{ji}^{t,m}
    :=
    (1-\rho_\nu)\sum_{\ell=0}^{m-1}\rho_\nu^{m-1-\ell}
    \bigl(d_j^{t+\ell}-d_i^{t+\ell}\bigr).
    \label{eq:window_gap}
\end{equation}
This quantity is fully observable from logged \(d\)-statistics and determines,
through \eqref{eq:m_step_z_exact}, how much separation is injected into the
\(\nu\)-gap over the window \([t,t+m-1]\).
\end{definition}

\begin{proposition}[Data-dependent \(\nu\)-separation certificate]
\label{prop:data_driven_nu_certificate}
Fix \(m\ge 1\). Suppose there exist \(T_0<\infty\) and \(\Gamma_m>0\) such that
for every \(t\ge T_0\),
\begin{equation}
    \min_{(j,i)\in\BB\times\HH}\mathcal G_{ji}^{t,m}\ge \Gamma_m.
    \label{eq:window_gap_lb}
\end{equation}
Then, for all \(t\ge T_0\),
\begin{equation}
    \min_{(j,i)\in\BB\times\HH} z_{ji}^{t+m}
    \ge
    \rho_\nu^m \min_{(j,i)\in\BB\times\HH} z_{ji}^t + \Gamma_m.
    \label{eq:block_gap_recursion}
\end{equation}
Consequently,
\begin{equation}
    \liminf_{k\to\infty}\;
    \min_{(j,i)\in\BB\times\HH} z_{ji}^{T_0+km}
    \ge
    \frac{\Gamma_m}{1-\rho_\nu^m}
    >
    0.
    \label{eq:block_gap_positive}
\end{equation}
\end{proposition}

\begin{proof}
By Lemma~\ref{lem:m_step_nu_gap},
\[
    z_{ji}^{t+m}=\rho_\nu^m z_{ji}^t+\mathcal G_{ji}^{t,m}.
\]
Taking the minimum over \((j,i)\in\BB\times\HH\) and using
\eqref{eq:window_gap_lb} yields \eqref{eq:block_gap_recursion}. Also,
let
\(
y_k:=\min_{(j,i)\in\BB\times\HH} z_{ji}^{T_0+km}.
\)
Then \(y_{k+1}\ge \rho_\nu^m y_k+\Gamma_m\). Unrolling this affine recursion
gives
\[
    y_k
    \ge
    \rho_\nu^{mk} y_0
    +
    \Gamma_m\sum_{s=0}^{k-1}\rho_\nu^{ms}.
\]
Since \(\sum_{s=0}^{k-1}\rho_\nu^{ms}\to (1-\rho_\nu^m)^{-1}\) and
\(\rho_\nu^{mk}y_0\to 0\), we obtain \eqref{eq:block_gap_positive}.
\end{proof}

\paragraph{Empirical support from the \(\rho_\nu=0.95\), threshold \(0.9\) run.}
The run recorded in
\texttt{hrac-log-label\_flipping0.95v0.9.txt} provides the clearest empirical
match to the group-level theory. In the late-stage window
\(t\in[18000,20000]\), the empirical mean gap
\(\bar\nu_B^t-\bar\nu_H^t\) stays positive, typically between \(0.06\) and
 \(0.15\) with median near \(0.12\), which is consistent with
Condition~\ref{cond:mean_nu_sep} and therefore supports
Theorem~\ref{thm:mean_nu_to_penalty}. Moreover, the best-fit curves for the
\(d\)-statistics are consistent with the premise of
Theorem~\ref{thm:mean_d_to_mean_nu}: the Byzantine group occupies a slightly
higher late-stage mean \(d\)-band, which the EMA then converts into a higher
late-stage mean \(\nu\)-band. The same run also shows threshold activation and
a clear weight split: threshold-crossing Byzantine clients repeatedly receive
weights around \(0.066\) to \(0.088\), whereas the honest majority remains near
\(0.108\) to \(0.111\). Thus the realized trajectory follows the intended pipeline
\[
    \text{positive late-stage bands}
    \;\Longrightarrow\;
    \text{higher Byzantine }\nu\text{-band}
    \;\Longrightarrow\;
    \text{threshold crossing}
    \;\Longrightarrow\;
    \text{weight suppression}.
\]
Figure~\ref{fig:hrac_iid_best_fit_d_nu} gives a paper-ready visualization of
this same run. It makes the intended interpretation explicit: neither
population is driven to a zero-statistics regime; instead, both populations
remain in positive late-stage bands, while the Byzantine trend is
systematically higher, especially for \(\nu\), which is precisely the statistic
used by the weighting rule.

\begin{figure}[t]
    \centering
    \includegraphics[width=\textwidth]{hrac-log-label_flipping0.95v0.9_d_nu_curves_with_best_fit_0_20000.png}
    \caption{Late-stage \(d\)- and \(\nu\)-statistics for the IID
    label-flipping run with \(\rho_\nu=0.95\) and activation threshold \(0.9\),
    together with per-client best-fit lines over \(t\in[0,20000]\). The upper
    panel plots the logged \(d_i^t\) trajectories, and the lower panel plots
    the corresponding \(\nu_i^t\) trajectories. Two features are visible.
    First, neither honest nor Byzantine clients collapse to \(0\); both groups
    remain in positive late-stage bands. Second, the Byzantine clients exhibit
    a higher fitted trend, especially in \(\nu\), than the honest clients.
    This supports the group-level mechanism formalized by
    Theorem~\ref{thm:mean_d_to_mean_nu} and
    Theorem~\ref{thm:mean_nu_to_penalty}: a larger late-stage Byzantine
    \(d\)-band is accumulated by the EMA into a larger late-stage
    \(\nu\)-band, which then activates stronger penalties and smaller
    aggregation weights.}
    \label{fig:hrac_iid_best_fit_d_nu}
\end{figure}

In the paper body, Figure~\ref{fig:hrac_iid_best_fit_d_nu} can therefore be
read as a compact visual summary of the adaptive-statistics story: the honest
and Byzantine clients both remain active in positive late-stage bands, but the
Byzantine group sits higher on average, and that relative elevation is exactly
what HRAC converts into down-weighting once the thresholded penalty mechanism
is activated. For clarity, the implementation used in this run applies the
threshold test \(\nu_i^t>0.9\) in the aggregation module; any remaining
\(\nu>1\) text in the log output should be read as a stale logger string
rather than the effective runtime threshold.

\begin{theorem}[Excess gap from late-stage \(\nu\)-band separation]
\label{thm:excess_from_d}
Assume Condition~\ref{cond:nu_sep_emp} holds.
Define the idealized excess score
\begin{equation}
    e_k^t := (\nu_k^t-\bar\nu_t)_+,
    \qquad
    \bar\nu_t := \frac1N\sum_{\ell=1}^N \nu_\ell^t.
    \label{eq:ideal_excess_score}
\end{equation}
Define
\[
    \Gamma_e
    :=
    \frac{H}{N}\gamma_\nu - \frac{H}{N}\omega_H - \frac{B}{N}\omega_B.
\]
If \(\Gamma_e>0\), then for every \(t\ge T_1\),
\begin{equation}
    \min_{j\in\BB} e_j^t - \max_{i\in\HH} e_i^t
    \ge
    \Gamma_e
    > 0.
    \label{eq:e_gap_ass}
\end{equation}
Hence with
\(E_H:=\sup_{t\ge T_1}\max_{i\in\HH}e_i^t\), we have
\begin{align}
    \max_{i\in\HH} e_i^t &\le E_H, \label{eq:e_h_bound}\\
    \min_{j\in\BB} e_j^t &\ge E_H+\gamma_e, \label{eq:e_b_bound}
\end{align}
for all \(t\ge T_1\), with any \(0<\gamma_e\le \Gamma_e\).
Moreover, with exact raw scores
\(\hat w_k^t=\frac{1}{1+\alpha e_k^t}\), the raw-weight ratio gap holds:
\begin{equation}
    \max_{j\in\BB}\hat w_j^t
    \le
    \kappa_t \min_{i\in\HH}\hat w_i^t,
    \qquad
    \kappa_t
    :=
    \frac{1+\alpha \max_{i\in\HH}e_i^t}
         {1+\alpha \min_{j\in\BB}e_j^t}
    \le
    \frac{1+\alpha E_H}{1+\alpha(E_H+\gamma_e)}
    =:\kappa<1.
    \label{eq:raw_ratio_gap}
\end{equation}
\end{theorem}

\begin{proof}
Fix \(t\ge T_1\), and denote
\[
    a_t:=\max_{i\in\HH}\nu_i^t,\quad
    \underline a_t:=\min_{i\in\HH}\nu_i^t,\quad
    b_t:=\min_{j\in\BB}\nu_j^t,\quad
    \overline b_t:=\max_{j\in\BB}\nu_j^t.
\]
By \eqref{eq:nu_h_spread}--\eqref{eq:nu_b_spread},
\(\underline a_t\ge a_t-\omega_H\), \(\overline b_t\le b_t+\omega_B\), and by
\eqref{eq:nu_gap_ass}, \(b_t-a_t\ge\gamma_\nu\).
Since
\[
    \bar\nu_t
    =
    \frac{1}{N}\!\left(\sum_{i\in\HH}\nu_i^t+\sum_{j\in\BB}\nu_j^t\right),
\]
we have
\[
    \bar\nu_t
    \ge
    \frac{H\underline a_t + B b_t}{N}
    \ge
    a_t - \frac{H}{N}\omega_H + \frac{B}{N}\gamma_\nu,
\]
and
\[
    \bar\nu_t
    \le
    \frac{H a_t + B\overline b_t}{N}
    \le
    b_t - \frac{H}{N}(b_t-a_t) + \frac{B}{N}\omega_B
    \le
    b_t - \frac{H}{N}\gamma_\nu + \frac{B}{N}\omega_B.
\]
Therefore
\[
    \max_{i\in\HH} e_i^t
    =
    \max_{i\in\HH}(\nu_i^t-\bar\nu_t)_+
    \le
    (a_t-\bar\nu_t)_+
    \le
    \left(\frac{H}{N}\omega_H-\frac{B}{N}\gamma_\nu\right)_+,
\]
and
\[
    \min_{j\in\BB} e_j^t
    \ge
    (b_t-\bar\nu_t)_+
    \ge
    \frac{H}{N}\gamma_\nu-\frac{B}{N}\omega_B.
\]
Subtracting yields \eqref{eq:e_gap_ass} with
\(\Gamma_e=\frac{H}{N}\gamma_\nu-\frac{H}{N}\omega_H-\frac{B}{N}\omega_B\).
Then \eqref{eq:e_h_bound}--\eqref{eq:e_b_bound} follow by definition of
\(E_H\) and any \(0<\gamma_e\le\Gamma_e\).

For the ratio bound, for any \(i\in\HH\), \(j\in\BB\),
\[
    \frac{\hat w_j^t}{\hat w_i^t}
    =
    \frac{1+\alpha e_i^t}{1+\alpha e_j^t}
    \le
    \frac{1+\alpha \max_{i'\in\HH}e_{i'}^t}
         {1+\alpha \min_{j'\in\BB}e_{j'}^t}
    =\kappa_t.
\]
Using \eqref{eq:e_gap_ass} and \(\max_{i\in\HH}e_i^t\le E_H\) gives
\(\kappa_t\le\frac{1+\alpha E_H}{1+\alpha(E_H+\gamma_e)}<1\).
\end{proof}

\begin{remark}[Relation to the implemented activation threshold]
\label{rem:ideal_excess_vs_threshold}
Theorem~\ref{thm:excess_from_d} is stated for the idealized positive-part score
\eqref{eq:ideal_excess_score}. The implementation further applies an activation
gate, penalizing only clients with \(\nu_i^t>0.9\). Therefore the theorem should
be read as the clean separation result for the underlying excess geometry. To
transfer it verbatim to the exact implemented score, one additionally needs a
late-stage threshold-alignment condition, for example that Byzantine clients
remain above the activation threshold while honest clients do not create a
competing excess band.
\end{remark}

\begin{proposition}[Threshold alignment transfers ideal excess separation to the implemented score]
\label{prop:threshold_alignment}
Assume Condition~\ref{cond:nu_sep_emp} and suppose the conclusions
\eqref{eq:e_h_bound}--\eqref{eq:e_b_bound} of
Theorem~\ref{thm:excess_from_d} hold for some \(E_H\ge 0\) and
\(\gamma_e>0\). Define the implemented score
\begin{equation}
    \xi_k^t
    :=
    \begin{cases}
        (\nu_k^t-\bar\nu_t)_+, & \nu_k^t>0.9,\\
        0, & \nu_k^t\le 0.9.
    \end{cases}
    \label{eq:implemented_xi}
\end{equation}
Suppose there exists \(T_\mathrm{act}<\infty\) such that for all \(t\ge
T_\mathrm{act}\),
\begin{equation}
    \min_{j\in\BB}\nu_j^t > 0.9,
    \label{eq:byz_threshold_on}
\end{equation}
and either of the following holds:
\begin{enumerate}[label=(\roman*)]
    \item \emph{Honest inactive regime:}
    \(
    \max_{i\in\HH}\nu_i^t \le 0.9
    \);
    \item \emph{Honest active-but-dominated regime:}
    \(
    \max_{i\in\HH} e_i^t \le E_H
    \)
    and
    \(
    \min_{j\in\BB} e_j^t \ge E_H+\gamma_e
    \)
    for some \(\gamma_e>0\).
\end{enumerate}
Then for all sufficiently large \(t\),
\begin{equation}
    \min_{j\in\BB}\xi_j^t - \max_{i\in\HH}\xi_i^t \ge \gamma_\xi > 0
    \label{eq:implemented_xi_gap}
\end{equation}
for some \(\gamma_\xi>0\). Consequently, with implemented raw scores
\(
\hat w_k^t = 1/(1+\alpha \xi_k^t),
\)
the Byzantine raw scores remain uniformly smaller than the honest raw scores.
\end{proposition}

\begin{proof}
Fix \(t\ge T_\mathrm{act}\). Under \eqref{eq:byz_threshold_on}, every Byzantine
client is active, so \(\xi_j^t=e_j^t\) for \(j\in\BB\). In case (i), all honest
clients are inactive, hence \(\xi_i^t=0\) for \(i\in\HH\), while
\eqref{eq:e_b_bound} gives \(\min_{j\in\BB}e_j^t\ge E_H+\gamma_e\). Therefore
\[
    \min_{j\in\BB}\xi_j^t-\max_{i\in\HH}\xi_i^t
    =
    \min_{j\in\BB}e_j^t
    \ge
    \gamma_e.
\]
In case (ii), active honest clients satisfy \(\xi_i^t=e_i^t\) and inactive
honest clients satisfy \(\xi_i^t=0\le e_i^t\), so
\[
    \max_{i\in\HH}\xi_i^t \le \max_{i\in\HH} e_i^t \le E_H,
\]
whereas \(\min_{j\in\BB}\xi_j^t=\min_{j\in\BB}e_j^t\ge E_H+\gamma_e\). Hence
\eqref{eq:implemented_xi_gap} holds with \(\gamma_\xi:=\gamma_e\). Monotonicity
of \(x\mapsto 1/(1+\alpha x)\) then transfers the score gap to the raw-weight
ordering.
\end{proof}

\begin{corollary}[Excess gap from a verified windowed \(d\)-certificate]
\label{cor:excess_from_window_certificate}
Suppose the hypotheses of Proposition~\ref{prop:data_driven_nu_certificate}
hold and the resulting positive \(\nu\)-gap level is large enough so that, for
all sufficiently large block times, the within-group spreads satisfy
\eqref{eq:nu_h_spread}--\eqref{eq:nu_b_spread} with
\[
    \gamma_\nu
    :=
    \frac{\Gamma_m}{1-\rho_\nu^m} > 0.
\]
Then the conclusions of Theorem~\ref{thm:excess_from_d} apply on those block
times. Hence the relative penalty mechanism admits a log-verifiable route from
windowed \(d\)-separation to excess-score separation.
\end{corollary}

\begin{proof}
Use Proposition~\ref{prop:data_driven_nu_certificate} to obtain a positive
block-time lower bound on \(\min_{j\in\BB}\nu_j^t-\max_{i\in\HH}\nu_i^t\), then
invoke Theorem~\ref{thm:excess_from_d}.
\end{proof}

\begin{theorem}[Weight control from the exact \(\nu\)-weight recursion]
\label{thm:weight_control_exact}
Suppose there exist \(T_0\), \(E_H\ge 0\), and \(\gamma_e>0\) such that
\eqref{eq:e_h_bound}--\eqref{eq:e_b_bound} hold for all \(t\ge T_0\), and that
the max-weight cap is inactive for \(t\ge T_0\), so
\begin{equation}
    w_i^t = \frac{\hat w_i^t}{\sum_{k=1}^N \hat w_k^t},
    \qquad
    \hat w_i^t = \frac{1}{1+\alpha e_i^t}.
    \label{eq:exact_weight_no_cap}
\end{equation}
Define
\[
    u_H^{\min} := \frac{1}{1+\alpha E_H},
    \qquad
    u_H^{\max} := 1,
    \qquad
    u_B^{\max} := \frac{1}{1+\alpha(E_H+\gamma_e)}.
\]
Then for all \(t\ge T_0\),
\begin{equation}
    \beta_t
    :=
    \sum_{j\in\BB} w_j^t
    \le
    \bar\beta
    :=
    \frac{B\,u_B^{\max}}{H\,u_H^{\min}+B\,u_B^{\max}}
    < 1.
    \label{eq:beta_bar_exact}
\end{equation}
Hence a concrete bound on \(\beta_t\) is verified.
\end{theorem}

\begin{proof}
From \eqref{eq:e_h_bound}--\eqref{eq:e_b_bound} and
\eqref{eq:exact_weight_no_cap}, honest and Byzantine unnormalised weights obey
\[
    \hat w_i^t \ge u_H^{\min}\quad (i\in\HH),
    \qquad
    \hat w_j^t \le u_B^{\max}\quad (j\in\BB).
\]
Therefore
\[
    \sum_{j\in\BB}\hat w_j^t \le B u_B^{\max},
    \qquad
    \sum_{i\in\HH}\hat w_i^t \ge H u_H^{\min}.
\]
After normalisation,
\[
    \beta_t
    =
    \frac{\sum_{j\in\BB}\hat w_j^t}
         {\sum_{i\in\HH}\hat w_i^t+\sum_{j\in\BB}\hat w_j^t}
    \le
    \frac{B\,u_B^{\max}}{H\,u_H^{\min}+B\,u_B^{\max}}
    =\bar\beta.
\]
\end{proof}

\subsection{Max-weight cap and iterative renormalisation}

In the current implementation, the inverse-excess raw scores
\(\hat w_i^t=1/(1+\alpha e_i^t)\) are first normalised to a probability vector,
and the max-weight cap is then applied to the \emph{normalised} weights:
\[
    w_i^t \gets \min\{w_i^t,w_{\max}\},
\]
followed by renormalisation. This cap--renormalise step is repeated only if
needed, for at most a fixed finite number of iterations in code.

\begin{lemma}[Inactive cap regime]
\label{lem:cap_inactive}
If the already normalised weights satisfy \(w_i^t\le w_{\max}\) for every
client \(i\) at round \(t\), then the cap step is inactive and the weights
coincide with the uncapped normalised \(\nu\)-score formula induced by
\eqref{eq:exact_weight_no_cap}.
\end{lemma}

\begin{proof}
Immediate.
\end{proof}

\begin{lemma}[Inactive cap when \(w_{\max}\ge 1\)]
\label{lem:cap_honest_safe}
For the inverse excess scores \(\hat w_i=1/(1+\alpha e_i)\) with \(e_i\ge 0\),
one always has \(\hat w_i\in(0,1]\). After normalisation, this implies
\(w_i\in(0,1]\) as well. Consequently, if \(w_{\max}\ge 1\), then
\(\min\{w_i,w_{\max}\}=w_i\) for every client \(i\), and a single
cap check leaves all weights unchanged.
\end{lemma}

\begin{proof}
Immediate from \(e_i\ge 0\Rightarrow 1+\alpha e_i\ge 1\).
\end{proof}

\begin{proposition}[Flattened invariant: cap-dominated weights]
\label{prop:cap_flat}
Suppose that after the finite cap--renormalise loop at round \(t\), the
returned weight vector is flat across clients. Then the final normalised
weights are uniform \(w_i^t=1/N\) and
\begin{equation}
    \beta_t = \frac{B}{N},
    \qquad
    \varepsilon_{w,t}
    =
    \sum_{i\in\HH}\left|\frac{1}{N}-\frac{1}{H}\right|
    =
    \frac{B}{N}.
    \label{eq:flat_beta_eps}
\end{equation}
\end{proposition}

\begin{proof}
If the returned weight vector is flat and sums to \(1\), every coordinate must
equal \(1/N\). The distortion identity is elementary algebra.
\end{proof}

\begin{proposition}[General active-cap fixed-point invariant (\(\beta_t\) control)]
\label{prop:beta_active_cap_general}
Consider any round \(t\) where the finite cap--renormalise loop returns a
weight vector \(w^t=(w_1^t,\dots,w_N^t)\) satisfying
\[
    w_i^t\ge 0,\qquad \sum_{i=1}^N w_i^t=1,\qquad w_i^t\le w_{\max}\ \forall i.
\]
Then, regardless of how many cap iterations were used within that loop,
\begin{equation}
    \beta_t
    :=
    \sum_{j\in\BB} w_j^t
    \le
    \min\{1,\;B\,w_{\max}\}.
    \label{eq:beta_active_cap_general}
\end{equation}
\end{proposition}

\begin{proof}
Since each Byzantine coordinate is bounded by \(w_{\max}\),
\[
    \beta_t=\sum_{j\in\BB}w_j^t\le \sum_{j\in\BB} w_{\max}=B\,w_{\max}.
\]
Also \(\beta_t\le 1\) because weights sum to \(1\). Combining gives
\eqref{eq:beta_active_cap_general}.
\end{proof}

\begin{remark}[Bridge from \(r_H\) to \(\omega_H\)]
\label{rem:rh_to_omegah}
Let
\(
\omega_H^t:=\max_{i\in\HH}\nu_i^t-\min_{i\in\HH}\nu_i^t
\).
Since each \(\nu_i^t\) is an EMA of \(d_i^t\), for any \(i,i'\in\HH\),
\[
    \nu_i^{t+1}-\nu_{i'}^{t+1}
    =
    \rho_\nu(\nu_i^t-\nu_{i'}^t)
    +(1-\rho_\nu)(d_i^t-d_{i'}^t).
\]
Hence if \(\max_{i\in\HH}d_i^t\le r_H\) eventually, then
\(\omega_H^t\) is uniformly bounded by \(O(r_H)\) after transients. We use this
as the explicit link between honest tracking smoothness and honest
\(\nu\)-spread control in the weight concentration proof. The conversion from
state-tube bounds such as \eqref{eq:honest_tube_eta} to a direct \(d\)-bound is
a separate implementation-geometry step.
\end{remark}

\begin{theorem}[Honest weight concentration from algorithmic smoothness]
\label{thm:honest_weight_concentration}
Suppose the honest \(d\)-statistics satisfy
\(\max_{i\in\HH} d_i^t \le r_H\) eventually. Then the intra-honest weight
distortion is structurally bounded by the algorithmic parameters and the
Byzantine mass \(\beta_t\):
\begin{equation}
    \varepsilon_{w,t}
    :=
    \sum_{i\in\HH}\left|w_i^t-\frac{1}{H}\right|
    \le
    C_H\,\alpha r_H + \beta_t,
    \label{eq:structural_epsw_bound}
\end{equation}
for some constant \(C_H>0\) depending only on a lower bound of the honest raw
weight level \(u_H^{-}\), where
\[
    u_H^{-}
    :=
    \inf_t \min_{i\in\HH}\hat w_i^t
    > 0.
\]
For sufficiently small learning rates \((\eta \to 0,\ r_H \to \mathcal{O}(\sigma^2))\),
the distortion simplifies to
\[
    \varepsilon_{w,t} \le C_w \beta_t + \mathcal{O}(\sigma^2),
\]
satisfying the strict refinement condition required by
Proposition~\ref{prop:tight_floor_beta2}.
\end{theorem}

\begin{proof}
By Remark~\ref{rem:rh_to_omegah}, since \(\max_{i\in\HH} d_i^t \le r_H\), the
EMA statistic \(\nu_i^t\) for all honest clients concentrates within a narrow
band \(\omega_H = O(r_H)\). Recall that the excess score used for weighting is
\[
    e_i^t = \max(0, \nu_i^t - \bar{\nu}_t).
\]
Because honest \(\nu_i^t\) are tightly clustered, their excess scores
\(e_i^t\) are equally clustered. Specifically,
\[
    \max_{i\in\HH} e_i^t - \min_{i\in\HH} e_i^t \le \omega_H = O(r_H).
\]
The raw weights are computed as
\[
    \hat{w}_i^t = \frac{1}{1+\alpha e_i^t}.
\]
Since \(f(x)=1/(1+\alpha x)\) has derivative magnitude at most \(\alpha\) on
\(x\ge 0\), the spread of the honest raw weights is bounded by
\[
    u_H^{+} - u_H^{-}
    =
    \max_{i\in\HH} \hat{w}_i^t - \min_{i\in\HH} \hat{w}_i^t
    \le
    \alpha \left( \max_{i\in\HH} e_i^t - \min_{i\in\HH} e_i^t \right)
    \le
    C_\omega\,\alpha r_H
\]
for some \(C_\omega>0\) from the \(\omega_H=O(r_H)\) relation. Using the
fundamental weight distortion inequality, the normalized weight deviation
\(\varepsilon_{w,t}\) satisfies
\[
    \varepsilon_{w,t}
    \le
    \frac{H(u_H^{+} - u_H^{-})}{\sum_{k=1}^N \hat{w}_k^t} + \beta_t.
\]
Since the sum of raw weights satisfies \(S_t \ge H u_H^{-}\), the first term is
bounded by
\[
    \frac{C_\omega\alpha r_H}{u_H^{-}}
    =
    C_H\alpha r_H
\]
with \(C_H:=C_\omega/u_H^{-}\). Therefore,
\[
    \varepsilon_{w,t} \le C_H\alpha r_H + \beta_t.
\]
\end{proof}

\begin{corollary}[Final model+\(\nu\) convergence statement]
\label{cor:final_model_nu}
Suppose Assumptions~\ref{ass:lower_bounded} to \ref{ass:honest_majority},
suppose moreover that one of the following late-stage routes holds:
\begin{enumerate}[label=(\alph*)]
    \item the \emph{strong pointwise route} of
    Remark~\ref{rem:pointwise_route_status}: either
    Condition~\ref{cond:nu_sep_emp} or the window-certificate route of
    Corollary~\ref{cor:excess_from_window_certificate}, together with the direct
    separation conclusion \eqref{eq:e_gap_ass}--\eqref{eq:raw_ratio_gap} in
    Theorem~\ref{thm:excess_from_d}, the inactive-cap regime of
    Theorem~\ref{thm:weight_control_exact}, and the honest-weight distortion
    regime of Theorem~\ref{thm:honest_weight_concentration} (or the symmetry
    case \(\varepsilon_{w,t}=0\));
    \item the \emph{group-level route}: Condition~\ref{cond:mean_nu_sep} and the
    resulting Byzantine penalty conclusion of
    Theorem~\ref{thm:mean_nu_to_penalty}, together with the exact
    implementation bridge supplied by
    Proposition~\ref{prop:threshold_alignment} once the activation regime is
    aligned on the realized trajectory.
\end{enumerate}
Then:
\begin{enumerate}[label=(\roman*)]
    \item \textbf{Model convergence:} The bound \eqref{eq:main_full} holds, so
    HRAC converges sublinearly to a neighbourhood of first-order stationarity.
    \item \textbf{Adaptive-state conclusion:} the strong late-stage route
    yields explicit excess-score separation and weight suppression through
    Theorem~\ref{thm:excess_from_d}, while the weaker group-level route yields
    Byzantine penalty pressure through Theorem~\ref{thm:mean_nu_to_penalty} and,
    once threshold alignment is realized, the corresponding suppression of
    Byzantine weights in the implemented rule.
    In addition, for each client \(i\), if
    \(d_i^t\to d_i^\star\), then \(\nu_i^t\to d_i^\star\) by
    Lemma~\ref{lem:nu_tracking}.
    More generally, if \(d_i^t\) remains in a bounded positive band, then
    \(\nu_i^t\) remains in the same asymptotic band by
    Lemma~\ref{lem:nu_tracking}.
\end{enumerate}
\end{corollary}

\begin{proof}
Item (i) is always given by Theorem~\ref{thm:main}. Under the strong pointwise
route, Theorem~\ref{thm:excess_from_d},
Theorem~\ref{thm:weight_control_exact}, and
Theorem~\ref{thm:honest_weight_concentration} make the weight-control inputs in
\eqref{eq:agg_error_final} concrete. Under the weaker group-level route, the
adaptive statistics still exhibit Byzantine penalty pressure and thresholded
weight suppression, although without the full explicit pointwise excess-gap
constants of the strong route.
Item (ii) is Lemma~\ref{lem:nu_tracking}.
\end{proof}

\section{Discussion}

The final bound \eqref{eq:main_full} shows that the convergence behaviour of
HRAC is controlled by five terms:
\begin{enumerate}[label=\arabic*.]
    \item stochastic variance of honest updates;
    \item client heterogeneity;
    \item global clipping bias on honest messages;
    \item residual clipping bias relative to each client's own history;
    \item honest-weight distortion and residual Byzantine mass.
\end{enumerate}

The main conceptual advantage of HRAC is visible directly in the theorem: the
method does not force all clients towards a single global centre. Instead, each
client is clipped only relative to its \emph{own history}. This makes the method
naturally more tolerant to non-IID honest behaviour, while still suppressing
adversarial influence through global norm capping and history-based
down-weighting.

The proof chain is modular. Theorem~\ref{thm:main} gives optimization
convergence once the aggregation floor is controlled; Lemma~\ref{lem:nu_tracking}
turns \(d\)-control into \(\nu\)-control; Theorem~\ref{thm:mean_d_to_mean_nu}
formalizes the experimentally best-matched route from a larger Byzantine
group-mean \(d\)-band to a larger group-mean \(\nu\)-band;
Theorem~\ref{thm:mean_nu_to_penalty} then converts that \(\nu\)-band elevation into
penalty pressure; and Proposition~\ref{prop:data_driven_nu_certificate} gives
an exact log-verifiable route from windowed \(d\)-gaps to \(\nu\)-separation.
The stronger pointwise route remains available through
Theorem~\ref{thm:nu_separation_main}, Theorem~\ref{thm:excess_from_d}, and
Proposition~\ref{prop:threshold_alignment}. As emphasized in
Remark~\ref{rem:adaptive_scope}, the cleanest theorem-level part of the
adaptive-statistics layer is the chain
\[
    \text{higher group-mean } d
    \;\Longrightarrow\;
    \text{higher group-mean } \nu
    \;\Longrightarrow\;
    \text{stronger penalty pressure},
\]
whereas the route from attack semantics to a higher Byzantine \(d\)-band is
presented as a code-aligned sufficient mechanism supported by the runs.

\paragraph{How to read the adaptive statistics.}
The experiments motivating this manuscript suggest that neither honest nor
Byzantine \(d\)-values converge to \(0\). Instead, both populations remain in
positive late-stage bands, while Byzantine clients exhibit a higher late-stage
\(\nu\)-band on average and in fitted trend lines. Accordingly, the main role
of the theory is to explain and formalize \emph{relative elevation} and its
weighting consequence, not to claim a universal zero-versus-nonzero
separation.

\paragraph{Conclusion.}
The manuscript now supports the following message cleanly: HRAC admits a
rigorous optimization bound, an exact \(d\to\nu\) recursion, and a
log-verifiable certificate route. On the adaptive-statistics side, the most
defensible theorem-level claim is the group-mean chain
\[
    \text{higher Byzantine } d\text{-band}
    \;\Longrightarrow\;
    \text{higher Byzantine } \nu\text{-band}
    \;\Longrightarrow\;
    \text{stronger penalty pressure and smaller weights}.
\]
The stronger pointwise separation route is retained as a sufficient refinement,
while the step from persistent attack mismatch to a higher Byzantine \(d\)-band
is stated as a code-aligned mechanism backed by the empirical runs. Among the
current experiments, the IID label-flipping configuration with
\(\rho_\nu=0.95\) and threshold \(0.9\) provides the clearest illustration of
this relative-elevation mechanism.
\end{comment}
\end{document}
